%% COVENANT -- anonymous supplement, cited by the journal and the dual-track versions ("Supplement Sec. X").
%% Section letters are fixed: A overview, B traces to the formal model, C capture, D UNDEFINED,
%% E checker snapshots and rounds, F per-engine triage, G dossier summary, H full results tables,
%% I the two fresh studies.
\documentclass[acmtog,anonymous,review]{acmart}

\usepackage{booktabs}
\usepackage{multirow}
\usepackage{xspace}
\usepackage{tikz}
\usepackage{xr}
\externaldocument[P-]{../journal/main}
\graphicspath{{../figures/}}
\setlength{\emergencystretch}{3em}
\pdfsuppressptexinfo=-1

\newcommand{\sys}{\textsc{Covenant}\xspace}
\newcommand{\rcgp}{RCGP\xspace}
\newcommand{\bench}{GPUBugBench\xspace}
\newcommand{\code}[1]{\texttt{#1}}
\newcommand{\xspec}{\ensuremath{\Rightarrow_{x}}}
\newcommand{\mspec}{\ensuremath{\Rightarrow_{m}}}
\newcommand{\spec}{\ensuremath{\Rightarrow_{\mathit{spec}}}}
\newcommand{\Haz}{\ensuremath{\mathit{Haz}}}
\newcommand{\posthoc}{\textsuperscript{\textsf{ph}}}
\newcommand{\relhash}{\code{f28a96b8b6fb}}
\newcommand{\prevhash}{\code{fa974bb2d72d}}
\DeclareRobustCommand{\circled}[1]{\tikz[baseline=(c.base)]\node[circle,fill=black!35,text=white,inner sep=.6pt,font=\sffamily\scriptsize\bfseries](c){#1};}

\setcopyright{none}
\acmJournal{TOG}
\acmVolume{0}
\acmNumber{0}
\acmArticle{0}
\acmYear{2026}
\acmMonth{0}
\acmSubmissionID{XXX}
\settopmatter{printacmref=false}
\renewcommand\footnotetextcopyrightpermission[1]{}

\begin{document}

\title{\textsc{Covenant}: Schedule-Independent Verification of Vulkan Frame Synchronization --- Supplementary Material}
\author{Anonymous Author(s)}

\begin{abstract}
This supplement accompanies the paper. It contains the translation from traces to events, with the treatment of
framebuffer-local dependencies and of pairs of layout transitions (\S B; the formal model, the algorithm and the
proofs are stated only in the paper); the design and evidence for host-write capture (\S C) and for spec-strict
transitions from \code{UNDEFINED} (\S D); the checker snapshots, what each study exposed, and the frozen blind rounds
(\S E); the per-study and per-engine triage (\S F); a summary of the 11 historical evidence dossiers (\S G); the full
results tables (\S H); and the protocols, triage rules and full results of the two fresh studies of the proved
configuration (\S I). Every number carries the checker snapshot and its provenance.
\end{abstract}

\maketitle
\hypersetup{pdfauthor={},pdfcreator={},pdfsubject={},pdfkeywords={},pdftitle={COVENANT: Supplementary Material}}

\appendix

%% ===================================================================
\section{Overview and Provenance}
\label{sup:overview}

\paragraph{How to read this supplement}
Section letters are fixed and are cited by the paper as ``Supplement~\S X''. Lettered sections point inside this
supplement; numbered definitions, lemmas, theorems, figures and appendices point to the paper, which holds the only
statement of the formal model and the proofs.

\paragraph{Provenance}
A result is \emph{fresh} when the checker snapshot was frozen, by source hash, before the study's first run, and the
result is used as scored. It is \emph{dev} (development data) when the checker was developed on that code or on
findings from it. It is \emph{post-hoc} when a later checker re-ran code already seen; post-hoc results are marked
\posthoc{} in tables. Fresh results are the primary evidence. The theorems of the paper describe one checker, the
release, in its verified configuration (Tables~\ref{tab:s-versions} and~\ref{tab:s-configs}). The other snapshots are
named by their role: the snapshot frozen for a study, the first snapshot of the verified configuration (frozen after
every blind round and every earlier engine study, so all its results are post-hoc), and the snapshots frozen before
each of the two fresh studies (\S I).

\begin{table}[htbp]
\caption{Checker snapshots. The hash is the SHA-1 over the sorted checker sources and the benchmark runner, as the
checker prints it.}
\label{tab:s-versions}
\scriptsize
\begin{tabular}{@{}p{0.27\columnwidth}p{0.15\columnwidth}p{0.48\columnwidth}@{}}
\toprule
Snapshot (role) & source hash & used for \\
\midrule
blind-round snapshots & per round & blind rounds 1--6, each frozen before its round (fresh) \\
tutorial-study snapshot & \code{8c006a42543b} & tutorial study and engine study 1 (fresh); early reference-semantics, barrier-mutant and ray-tracing/mesh tests (dev) \\
engine-study-2 snapshot & \code{d13c9184ac04} & engine study 2 (fresh); whole-frame differential test after its fixes (dev) \\
engine-study-3 snapshot & \code{2f74a46b1341} & engine study 3 (fresh); the dossier counts \\
first snapshot of the verified configuration & \code{45f2ce1a5046} & post-hoc re-runs of the earlier corpora \\
fresh-study-1 snapshot & \code{7cb216ed5753} & fresh study 1 (32 programs); barrier mutants on the pinned validation commit (post-hoc) \\
fresh-study-2 snapshot & \code{76e393826ae9} & fresh study 2 (25 programs); the source-patch cases \\
release predecessor & \prevhash & all post-hoc re-runs in the paper's release table; the regenerated whole-frame test \\
\textbf{the release} & \relhash & the configuration the paper's theorems describe; differs from its predecessor only in reading specialization constants, with the same access classification on every corpus shader module, so the re-runs above are its results \\
\bottomrule
\end{tabular}
\end{table}

\begin{table}[htbp]
\caption{Tested software. A seven- to twelve-character entry is the start of a git commit ID in the project's public
repository, so the exact code can be checked out; other entries are release versions.}
\label{tab:s-software}
\scriptsize
\begin{tabular}{@{}p{0.36\columnwidth}p{0.2\columnwidth}p{0.36\columnwidth}@{}}
\toprule
Project & version & used in \\
\midrule
Khronos validation layer & \code{0ad67fd} & baseline of every comparison \\
lavapipe (Mesa) & 25.2.8 & driver of every validation run \\
Vulkan specification (Vulkan-Docs) & v1.4.365 & cited text (line numbers from \code{193ebe195f47}) \\
Vulkan-Tutorial & \code{baf622df} & tutorial study \\
vulkan-guide & \code{dcf72a8b} & tutorial study \\
vkQuake & \code{749b4fd} & engine study 1 \\
Diligent Engine & \code{efe5da2} & engine study 1 \\
SDL3 & \code{e4af6f8} & engine study 2 \\
bgfx & \code{73183be} & engine study 2 \\
Granite & \code{fb178c8} & engine study 3 \\
VulkanSceneGraph & \code{2feabc6} & engine study 3 \\
vsgExamples & \code{46a82943} & engine study 3 \\
SaschaWillems/Vulkan & \code{41a4410} & development data \\
KhronosGroup/Vulkan-Samples & \code{177edebf} & development data \\
vuk & \code{61abde9} & development data \\
Qt & 6.4.2 & fresh study 1 \\
vkcube & 1.3.275 & fresh study 1 \\
GTK & 4.14.5 & fresh study 2 \\
vkmark & \code{ab6e6f3} & fresh study 2 (the tested Ubuntu package is source-identical) \\
mpv / libplacebo & 0.37 / 6.338 & fresh study 2 \\
GStreamer & 1.24.2 & fresh study 2 \\
\rcgp & \code{85f2b04} & expressiveness comparison (\S\ref{sup:rcgp}) \\
\bottomrule
\end{tabular}
\end{table}

\begin{table}[htbp]
\caption{Configurations of the checker. Only the verified configuration, the default, is covered by the theorems.
The other switches are opt-in and can be combined; each is recorded in the report's metadata. ``Unrecorded subpass
flags'': how a trace from a driver that did not record subpass-dependency flags is read.}
\label{tab:s-configs}
\small
\resizebox{\columnwidth}{!}{%
\begin{tabular}{@{}llllll@{}}
\toprule
Configuration & from \code{UNDEFINED} & host-write capture & unrecorded subpass flags & slot accesses & theorems \\
\midrule
\textbf{verified} (default) & strict & conservative (pages) & local & plain & \textbf{yes} \\
observed capture & strict & observed (bytes) & local & plain & no \\
heuristic \code{UNDEFINED} relaxation (R1, opt-in) & relaxed & conservative & local & plain & no \\
legacy global subpass flags (opt-in) & strict & conservative & global & plain & no \\
heuristic slot exemption (opt-in) & strict & conservative & local & exempt & no \\
\bottomrule
\end{tabular}}
\end{table}
\begin{sloppypar}
The switches are the environment variables \code{COVENANT\_HOST\_CAPTURE=observed},
\code{COVENANT\_HEURISTIC\_DISCARD=1}, \code{COVENANT\_LEGACY\_SUBPASS\_FLAGS=global} and
\code{COVENANT\_HEURISTIC\_SLOTS=1}; explicit arguments of the checking function override them, and the report
states the configuration the check actually ran with. The snapshots of
Table~\ref{tab:s-versions} other than the release read no subpass-dependency flags at all (\S\ref{sup:fblocal}); for
them, ``verified configuration'' means strict transitions and conservative capture.
\end{sloppypar}

%% ===================================================================
\section{From Traces to the Formal Model}
\label{sup:formal}

The formal model, the verification algorithm, the theorems and every proof are stated once, in the paper (Sections~4
and~5, Section~8 and Appendix~A), and are not repeated here, so that the submission has a single statement of its
guarantees. In particular, the paper makes no claim about reusing a barrier plan beyond the recorded execution, and a
repair is accepted only when the whole frame, compared by racing event pairs, reports no pair that the original frame
did not (paper, Section~5). This section adds the translation from the trace to events, how framebuffer-local
dependencies are read, the matrix of transition pairs, and the lesson of the barrier optimizer. Numbered references
to definitions, lemmas, theorems and figures point to the paper.

\subsection{From Trace to Events}
\label{sec:impl}
The model driver \code{libcovvk} is a stand-in \code{libvulkan} that returns fake handles and records a JSON-lines
trace: object creation (with formats and usages), descriptor writes and command streams, submissions with their
semaphores and fences, swapchain acquire and present, and host writes to mapped memory (\S\ref{sec:capture}). Shaders
are lifted from SPIR-V, and the C++ host is read with libclang for struct layouts.
\begin{itemize}
\item \textbf{Submissions.} A submission is preceded by a host signal that the queue waits on. Then come each
  batch's waits, the commands of its command buffers as recorded at submit time, its signals, and the fence.
\item \textbf{Draws.} Draws read their bound descriptors and their vertex, index and indirect buffers, and write
  their attachments. Byte ranges come from descriptor ranges, dynamic offsets and index ranges.
\item \textbf{Shader accesses.} A per-entry-point analysis of the SPIR-V call graph yields each shader's reads and
  writes and classifies storage bindings as atomic-only or plain; an atomic counts only at a scope that holds
  every invocation of the device, and an access through an index taken from an atomic is a plain access (exempting
  such ``slot'' accesses is an opt-in heuristic, \code{COVENANT\_HEURISTIC\_SLOTS=1}, outside the verified
  configuration). Atomic-only is decided per byte range of the bound buffer (a block member's offset plus the
  descriptor's offset), not per member name. A pointer access chain (\code{OpPtrAccessChain}) first steps by whole
  elements of its base type; only an \code{OpConstant} Element of 0 keeps the access inside the member it then names (a specialization constant may be overridden),
  otherwise the access may touch the whole range. (A nonzero Element on a pointer directly to a block is not valid
  SPIR-V, so a nonzero Element occurs only on pointers into arrays, where the conservative range is sound.) A matrix's size is bounded for both column- and row-major
  layouts. A pointer keeps every descriptor it may come from, through helper calls,
  returns, \code{OpPhi}, \code{OpSelect} and pointers held in variables. When a pipeline's shader code is missing or
  cannot be parsed, or an access reaches descriptor memory through a value the analysis cannot trace to a descriptor,
  every storage descriptor it binds is taken as written at every stage, and the gap is recorded (``shader code not
  available'' or ``shader accesses not attributed''), so the check is not guarantee-eligible. The same gap is
  recorded when an instruction outside the inventory of memory effects (Table~\ref{tab:s-effects}) takes a
  descriptor pointer or handle as an operand. A descriptor access is
  charged at the stages whose entry point uses the binding; the set layout's stage flags are only the fallback when
  no entry point is known. (The snapshots up to engine study 3 charged every stage the layout granted, which is sound
  but caused a false alarm on \bench case 94.)
\item \textbf{Descriptor arrays.} A direct draw's read of a sampled-image array is narrowed to the elements its index
  can select when the index is a flat 32-bit vertex attribute, passed through at most four shift/mask operations,
  whose bytes the host wrote before the submission, that no device command writes and that no host write races
  (paper, \S\ref{P-sec:guards}); otherwise every written element is read. The guard is re-decided whenever a frame is
  re-checked. The snapshots of both fresh studies read every written element in all cases, which is sound and is the
  false-alarm cause GTK-DESC of study~2 (\S I).
\item \textbf{Semaphores.} Each signal of a binary semaphore is a separate instance. A timeline wait is matched first
  against the initial value, then against the first sufficient signal.
\item \textbf{Swapchain.} \code{vkAcquireNextImageKHR} places three events on $\mathit{present}{:}i$: a wait for the
  previous present of image $i$, a read of the image, and the acquire signal. \code{vkQueuePresentKHR} places the
  present-semaphore wait, a read in the present layout and a ``presented'' signal on the same presentation queue.
\item \textbf{Compute, mesh shading and ray tracing.} Dispatches and mesh-task draws access their descriptors at the
  compute, task and mesh stages; indirect and count buffers are read at \code{DRAW\_INDIRECT}. An acceleration
  structure is the byte range of its backing buffer. A build writes its destination, reads and writes its scratch
  range, and reads its inputs and, for a top-level build, every bottom-level structure its instances reference
  (A-RT1). \code{vkCmdTraceRaysKHR} reads its shader-binding-table regions and its descriptors at the ray-tracing
  stage.
\item \textbf{Stage order.} Logically earlier and later stages follow the synchronization pipelines of the Vulkan
  registry. Expansion applies to execution scopes only.
\item \textbf{Events and ownership.} \code{vkCmdWaitEvents} is a barrier whose first scope contains only what precedes
  the latest \code{vkCmdSetEvent} of that event on the queue (A-EV1). Queue-family ownership is tracked per byte
  range, and per subresource for images.
\item \textbf{Unmodelled commands.} The driver records every command it does not implement with its command buffer
  and position, and the checker emits \code{timing.unmodelled-command} when that buffer is submitted. The verdict
  makes no claim for such a command.
\end{itemize}
\begin{table*}[t]
\caption{The shader analysis's inventory of memory effects: the class of every SPIR-V instruction that may take a
descriptor pointer or handle (image, sampled image, acceleration structure) as an operand. Any other instruction on
one makes the check not guarantee-eligible (model gap ``shader accesses not attributed''), and every storage
descriptor of the pipeline is then taken as written.}
\label{tab:s-effects}
\scriptsize
\begin{tabular}{@{}p{0.08\textwidth}p{0.88\textwidth}@{}}
\toprule
Effect & Instructions \\
\midrule
write & \code{OpStore}; every \code{OpAtomic*} except \code{OpAtomicLoad} (image atomics through
\code{OpImageTexelPointer}); the target of \code{OpCopyMemory} and \code{OpCopyMemorySized}; \code{OpImageWrite};
\code{OpCooperativeMatrixStore\{KHR,NV,TensorNV\}}, \code{OpCooperativeVectorStoreNV},
\code{OpCooperativeVector\{ReduceSum,OuterProduct\}AccumulateNV}; a pointer operand of a \code{GLSL.std.450}
instruction (\code{Modf}, \code{Frexp}) \\
read & \code{OpLoad}, \code{OpAtomicLoad}, \code{OpArrayLength}; the source of a copy; \code{OpImageRead},
\code{OpImageFetch}, \code{OpImage*Gather}, \code{OpImageSample*}, \code{OpImageSparse*}, \code{OpImageQuery*};
cooperative-matrix and cooperative-vector loads and multiplies; \code{OpTraceRayKHR}, \code{OpTraceNV},
\code{OpTraceRayMotionNV}, \code{OpRayQueryInitializeKHR}, \code{OpHitObjectTraceRay*NV} \\
passes on & access chains, \code{OpCopyObject}, \code{OpBitcast}, \code{OpImageTexelPointer},
\code{OpSampledImage}, \code{OpImage}, \code{OpPhi}, \code{OpSelect}, \code{OpCopyLogical}, \code{OpFunctionCall},
\code{OpReturnValue}, pointer comparisons; the result keeps every descriptor of its operands \\
none & \code{OpVariable}, \code{NonSemantic} debug instructions, control flow \\
\bottomrule
\end{tabular}
\end{table*}

\emph{Structure checks} (secondary, no guarantee) use the same trace: stage interfaces, host/shader struct layouts,
bindings used but not bound, push-constant ranges and framebuffer extents.

\subsection{Framebuffer-local dependencies}
\label{sup:fblocal}
The model driver records the \code{dependencyFlags} of every subpass dependency and barrier. A link with
\code{BY\_REGION}, or \code{VIEW\_LOCAL} with view offset 0, keeps its global meaning for a stage pair that involves a
stage outside framebuffer space; for two framebuffer-space accesses it orders only attachment accesses at the same
position (the same image view base; a single-sample input-attachment read counts); storage, sampled and uniform
accesses are never ordered by it. A non-zero view offset orders no framebuffer-space pair. A chain with a local link is
local (paper, Definition~\ref{P-def:spec} and \S\ref{P-sec:fblocal}). \code{VIEW\_LOCAL} is treated like
\code{BY\_REGION}, which is stricter than the specification. A trace from a driver that did not record subpass flags is
read as local, the weaker reading, with a note in the report; the opt-in legacy switch reads such flags as global and
leaves the verified configuration (Table~\ref{tab:s-configs}).

The snapshots frozen for the earlier studies, including both fresh studies, read no dependency flags and gave every
dependency its global meaning. On a frame with a \code{BY\_REGION} or \code{VIEW\_LOCAL} dependency they can therefore
accept a pair that the release reports. The
release's results on the traces of both fresh studies are in \S\ref{sup:release-posthoc}.

\subsection{Transition pairs}
\label{sup:ltpairs}
Two layout transitions of one image are exempt from the race check (exemption E5 of Definition~\ref{P-def:spec}, rule
R-LT) only when they are on one queue and at least one of them is performed by an image memory barrier, because only
then does the specification order them. Table~\ref{tab:s-ltmatrix} gives every kind of pair, the specification text
that decides it, and how the checker treats it. Ordering among transitions adds no chain: ``an application must not
rely on ordering of image layout transitions to influence ordering of other commands''.

\begin{table}[h]
\caption{Pairs of layout transitions on one image with no dependency chain between them (synchronization.adoc
262--269, Vulkan-Docs \code{v1.4.365}). ``Automatic'': a transition performed by a render pass; ``explicit'': one
performed by an image memory barrier.}
\label{tab:s-ltmatrix}
\scriptsize
\begin{tabular}{@{}p{2.9cm}p{2.5cm}p{1.1cm}l@{}}
\toprule
pair & specification & fresh-study snapshots & release \\
\midrule
explicit / explicit, one queue & ``in submission order'' & exempt & exempt \\
explicit / automatic, one queue & ``including those performed by render passes'' & exempt & exempt \\
automatic / automatic, one queue & not covered & exempt & \textbf{checked} \\
any pair on two queues & ``submitted to the same queue'' & checked & checked \\
transition / other command & ``must not rely on'' & no edge & no edge \\
\bottomrule
\end{tabular}
\end{table}
Each row has a unit test, and the third row a litmus program (two render passes with and without explicit external
dependencies). Checking automatic/automatic pairs reports new pairs in four GTK and two Qt programs of the fresh
studies and in \bench case 45; no diagnostic site is lost. Whether two automatic transitions are in fact ordered is
one of the questions the specification leaves open (\S\ref{sup:vsg5}); the checker reports them, which is the report
direction.

\subsection{Verified Barrier Transformations}
\label{sec:elide}
Because the checker decides hazard-freedom for the whole recorded frame, it can serve as the oracle of a barrier
optimizer. A greedy search visits barrier calls in recording order and, for each, tries removal and then narrower
masks built from the stages and flags of the accesses around each occurrence, keeping a candidate when the whole
frame, with all changes accepted so far, reports no hazard beyond those of the original frame. Barriers with
\code{HOST} stages, ownership transfers, masks the trace could not name and the barriers around unmodelled commands are
pinned; layout transitions are only narrowed. The release compares the frames by racing event pairs, matched by
identity, and re-verifies both plans pairwise. The measurements of \S\ref{sec:eval-elide} were made with an early
optimizer and checker (development data), which compared diagnostic keys rather than event pairs; they were not
re-run with the release.

\paragraph{The soundness lesson}
Running the early optimizer on eight unmodified samples exposed two gaps between its implementation and the theorem's
premise, both fixed since. \emph{Existing hazards mask new ones:} on a frame that already races, ``no new
diagnostic'' is not monotone, because one existing witness can hide new races. Without a guard the optimizer narrowed
all 93 per-frame transitions of one Khronos sample, and that plan creates 1,560 new races under the reference
semantics; the optimizer now pins every barrier between the two accesses of an existing hazard. \emph{The
presentation engine's read before the first acquire:} the strengthened frame charged that read only from an image's
first present on, so the optimizer narrowed an acquire transition's source stage below the semaphore's wait stage, and
validation reports the resulting write-after-read hazard. The strengthened frame now places a presentation-engine read
before every acquire signal.


%% ===================================================================
\section{Host-Write Capture}
\label{sec:capture}

Host writes to mapped memory are the part of the trace that a recording driver does not see as calls, so they need
their own argument: premise A2-H of Table~\ref{P-tab:assume}. This section gives the mechanism, the conservativity
argument, its exact preconditions, the tests, and the measured precision cost of the conservative default. Every
measurement here was made with the first snapshot of the verified configuration and is post-hoc.

\subsection{Mechanism}
The model driver allocates every mappable \code{VkDeviceMemory} as an anonymous private mapping and, at
\code{vkMapMemory}, registers it as a \emph{trap region} and makes it read-only. The first store to a page of a
read-only region raises \code{SIGSEGV}; the driver's handler marks that page dirty in the region's bitmap, records the
faulting thread (the first four distinct ones per interval, then an overflow flag), makes \emph{that page} writable and
returns, so the store is retried and succeeds. Later stores to the same page in the same interval do not fault.

Every Vulkan entry point of the driver takes the driver lock and, before anything else, flushes. For each region with
dirty pages the flush (1) collects the dirty page offsets and clears their bits; (2) re-protects the whole region
read-only \emph{before} looking at its bytes, so that a store racing the flush either lands before the re-protect and
is in this record, or faults and is in the next one; and (3) emits one host-write record with the memory object, the
dirty pages, the page size, the faulting threads, the threads of the previous and the current call, the number of live
threads, and the byte runs that differ from the driver's shadow copy of the mapping. Only then is the call's own record
written. \code{vkFreeMemory} unregisters the region; the driver's own writes into a mapping (the read-back of a
device-to-host copy) update the shadow and do not mark pages. The trace ends with a final flush.

\subsection{Capture intervals and event placement}
Let $C_1, C_2, \dots$ be the calls in the order the driver lock serializes them. The capture interval $I_k$ is the time
between the end of $C_{k-1}$'s flush and the start of $C_k$'s. Every host store to a trapped page during $I_k$ is in the
record emitted at the start of $C_k$: the page is read-only at the start of $I_k$, because the previous flush
re-protected it and only a fault re-opens a page, so the first store faults, and a later store hits a page already
marked. The record becomes host write events (stage \code{HOST}, access \code{HOST\_WRITE}) placed on the host queue
after $C_{k-1}$ and before $C_k$.

Every operation that orders host work against device work is a call on the host queue: submission (host to device),
fence and semaphore waits, completion queries that observe completion, queue and device idle waits, and acquire
(device to host). Each is some $C_j$, so it lies wholly before or after any interval. Placing the write inside $I_k$
puts it on the same side of every such call as the actual store, provided the store is ordered with the calls in trace
order (P3 below).

\subsection{What is charged}
\emph{Conservative} capture, the default: for each dirty page, every buffer or image bound to the memory object
that overlaps the page receives a host write of its bytes inside the page (an image: its whole subresource set).
Byte contents are not consulted. A binding made by the very next call (the bump-allocator pattern) is included.
\emph{Observed} capture (opt-in; the only mode of the snapshots frozen for the tutorial study and the three engine
studies): only byte runs that differ from the shadow are charged (runs closer
than 16 bytes merged; past the run budget, whole pages from the truncation point), and a page shared by a fresh
allocation and an older neighbour is charged to the fresh one only.

\subsection{Conservativity argument}
Let $\mathit{St}$ be the host stores to mapped device memory and $\mathit{Ev}$ the host write events produced. The
theorem needs: for every store $s \in \mathit{St}$ to byte $b$ of a binding $R$, an event $w \in \mathit{Ev}$ that writes
$b$ of $R$ and lies on the host queue between the same two calls as $s$, so that it is related to every device event as
$s$ is. Given P1--P4:
\begin{itemize}
\item \emph{Seen.} $s$ happens in some interval $I_k$ (P3). Its page $p$ is read-only at the start of $I_k$. If $s$ is
  the first store to $p$ in $I_k$, it faults and marks $p$ (P1: the handler receives the fault; P2: the store goes
  through the trapped mapping); otherwise an earlier store marked $p$. So $p$ is in the record of $C_k$.
\item \emph{Charged.} In conservative mode every byte of every binding overlapping $p$ is written by an event of that
  record, so $b$ of $R$ is (P4).
\item \emph{Placed.} The record's events lie between $C_{k-1}$ and $C_k$ on the host queue, as $s$ does.
\end{itemize}
Hence $\mathit{Ev}$ over-approximates $\mathit{St}$ in content (whole pages, all overlapping bindings) and is exact in
position up to the interval, which is all that host happens-before can distinguish. A same-value store, and a store
restored within the interval, fault and are therefore accesses. \emph{In observed mode the argument fails at
``Charged''}: a store that leaves the bytes unchanged produces no byte run and no event; that mode is outside the
theorem.

Not covered: host \emph{reads} of mapped memory are not trapped and are not events, so a device write racing a host
read is not reported; kernel stores on the application's behalf (\code{read(2)} into a mapping) fail with
\code{EFAULT} while the page is protected, which changes behaviour rather than escaping silently; memory mapped through
a path the driver does not see is outside P2.

\subsection{Preconditions}
\begin{itemize}
\item \textbf{P1 (trap installed).} The driver's handler is in effect and the mapping is a trap region. Checked: the
  driver re-reads the handler at every flush and records a capture warning if the application replaced it; a mapping
  that could not be trapped is recorded as untracked. Both become \code{timing.unmodelled} warnings, so a clean
  verdict says that it does not cover those writes.
\item \textbf{P2 (stores through the mapping).} Stores go through the pointer \code{vkMapMemory} returned, or pointers
  derived from it. Not checkable from the trace.
\item \textbf{P3 (one ordering thread).} Each store is ordered, in host happens-before, between the two calls that
  bound its interval in trace order. It holds when stores and calls are made by one thread. Checked as far as the trap
  allows: a record whose faulting thread differs from the thread of either bounding call, or whose thread list
  overflowed, is reported (``host writes from a thread other than the one making the surrounding Vulkan calls'').
  Records made while the process had more than one thread are counted but not reported, because only the first store
  per page faults.
\item \textbf{P4 (binding known).} The bytes belong to a buffer or image bound at that offset by the end of the
  interval, or by the call that ends it.
\end{itemize}

\subsection{Tests}
Unit tests with synthetic records check both modes: the default is conservative; a page-only record is charged in
both modes; a record with an empty byte diff is charged in conservative mode only; a store into the neighbouring buffer
of the same page is charged to both buffers in conservative mode and to the neighbour only in observed mode; a
cross-thread record and an untracked mapping are reported. An end-to-end probe on the model driver uses one coherent
host-visible allocation holding two 256-byte buffers, \code{src} and \code{nbr}, in one page; a copy reads \code{src},
and the store under test is made after the submit and before the fence wait (Table~\ref{tab:s-probe}). A trace test
checks that the racing record of the \code{same} and \code{restore} probes has an empty byte diff and one dirty page.

\begin{table}[htbp]
\caption{Host-capture probes on the model driver.}
\label{tab:s-probe}
\small
\resizebox{\columnwidth}{!}{%
\begin{tabular}{@{}llccc@{}}
\toprule
Probe & what the store does & race? & conservative & observed \\
\midrule
\code{ok} & no store before the wait & no & clean & clean \\
\code{changed} & a different value into \code{src} & yes & reported & reported \\
\code{same} & the value \code{src} already holds & yes & \textbf{reported} & missed \\
\code{restore} & a different value, then the original back & yes & \textbf{reported} & missed \\
\code{neighbour} & into \code{nbr} only (same page) & no & reported (page cost) & clean \\
\code{thread} & \code{changed}, from a second thread & yes & reported + P3 warning & reported \\
\bottomrule
\end{tabular}}
\end{table}

\subsection{Precision cost of the conservative default}
Table~\ref{tab:s-capcost} compares conservative with observed capture on the same trace files, both with the first
snapshot of the verified configuration; everything else is equal. Raw counts are timing diagnostics; groups are timing report groups.

\begin{table*}[t]
\caption{Conservative versus observed capture, first snapshot of the verified configuration\posthoc, same traces. GPUBugBench status: detected / fixed
variant not clean / expected miss / missed.}
\label{tab:s-capcost}
\small
\begin{tabular}{@{}p{0.2\textwidth}p{0.13\textwidth}p{0.13\textwidth}p{0.48\textwidth}@{}}
\toprule
Corpus & observed & conservative & what the difference is \\
\midrule
GPUBugBench, 114 cases & 81 / 1 / 2 / 30 & 78 / 4 / 2 / 30 & \textbf{recall gain scored as a loss}: 65, 66 and 68 become fixed-not-clean and 69 gains a code on both variants. All four are a true same-value race in the shared harness (two frames in flight, one uniform buffer, the host rewrites the unchanged camera after waiting only its own frame's fence). Blind timing columns unchanged (29 / 33 of 35) \\
Whole-frame, 1,251 programs & 97.9\% / 95.4\% & 97.9\% / 95.4\% & no host write races in a way the modes distinguish \\
RT / mesh, 15 modes & 15 / 15 & 15 / 15 & none \\
Stored samples, 28 (earlier driver) & 30,244 raw, 113 / 7 groups & same & earlier traces have no page-only records \\
Sascha T1--T6, newly recorded traces & 502 host-race raw & 1,607 & \textbf{recall gain}, all same-value stores: T2 0$\to$8, T3 0$\to$4, T4 0$\to$9, T5 0$\to$576, T6 2$\to$10, T1 500$\to$1,000 (the view uniform buffer) \\
Tutorials, 43 traces & 3,291 raw / 22 groups & 9,606 / 58 & \textbf{precision cost} (+6,315 raw / +36 groups; +24 scored false-alarm groups), all in vulkan-guide chapters 5 and 6: the per-frame scene buffer is a fresh 240-byte allocation in a page that also holds the in-flight frames' scene buffers and the material constants \\
SDL3 + bgfx, 12 & 221 / 21 / 0 & 221 / 21 / 0 & none \\
Granite + VSG, 12 & 1,060 / 13 / 9 & 1,060 / 13 / 9 & one P3 warning per Granite program in both modes \\
Diligent GLTFViewer & 126 / 4 / 2 & 136 / 14 / 2 & \textbf{precision cost}, +10 groups: the staging ring sub-allocates each upload next to earlier uploads a copy may still read; all page neighbours \\
vkQuake, 3 timedemos & 71,104 raw / 14 groups & 81,161 / 27 & \textbf{precision cost}, +13 groups, all false alarms: the lightstyle uniform buffer is double-buffered by dynamic offset and every changed byte lies in the half the GPU does not read (4); the lights buffer is a page neighbour (3); the two halves of the visibility buffer, selected by a push-constant offset, share a page (6). Of the added raw, 9,800 are same-value records, most on the bmodel ring, joining existing false-alarm groups \\
\bottomrule
\end{tabular}
\end{table*}

\paragraph{Reading}
Conservative capture is what the theorem needs. On these corpora it buys the same-value detections (Sascha T2--T5,
more of T6 and T1, and the true harness race in \bench 65, 66, 68 and 69). Its cost is page-granularity false alarms in
programs that sub-allocate host-visible memory below page size while earlier sub-allocations are in flight. vkQuake and
Diligent fall from 7 true / 2 technicalities / 11 false alarms of 20 groups to 7 / 2 / 34 of 43. Observed capture
reproduces the host-write results of the engine-study-3 snapshot exactly, but is outside the theorem.

\subsection{Record counts}
Table~\ref{tab:s-records} counts how often the observed mode would have been blind (records whose page faulted but
whose byte diff is empty) and how often P3 is at stake. Traces made by the earlier driver of the sample and tutorial
studies have no page-only records and no thread fields.

\begin{table}[htbp]
\caption{Host-write records, conservative mode, loaded by the first snapshot of the verified configuration. Driver:
the earlier driver of the sample and tutorial studies, or the current one.}
\label{tab:s-records}
\scriptsize
\resizebox{\columnwidth}{!}{%
\begin{tabular}{@{}lrrrrr@{}}
\toprule
Traces (driver) & traces & records & same-value & cross-thread (reported) & while $>1$ thread \\
\midrule
stored Sascha (earlier) & 22 & 1,234 & 0 & 0 & 0 \\
stored Khronos (earlier) & 6 & 152 & 0 & 0 & 0 \\
tutorials (earlier) & 43 & 907 & 79 & 0 & 0 \\
Sascha T1--T6 (current) & 6 & 191 & 88 & 0 & 0 \\
SDL3 (current) & 6 & 715 & 231 & 0 & 0 \\
bgfx (current) & 6 & 12,151 & 1,145 & 0 & 12,151 \\
Granite (current) & 4 & 3,205 & 840 & 182 & 3,205 \\
VSG (current) & 8 & 105 & 26 & 0 & 0 \\
Diligent (current) & 1 & 94 & 72 & 0 & 94 \\
vkQuake (current) & 3 & 14,645 & 5,304 & 4,964 & 14,645 \\
\bottomrule
\end{tabular}}
\end{table}

Same-value records are frequent wherever a program rewrites constants every frame; they matter only when the buffer is
in flight, the T1--T6 pattern. Cross-thread records occur only in Granite and vkQuake, which fill mapped memory from
worker threads; each of their programs carries one P3 warning, so its clean pairs are outside P3.

%% ===================================================================
\section{Transitions from \texttt{UNDEFINED}}
\label{sup:undefined}

\subsection{Two readings, and what the specification says}
A write followed by an image layout transition from \code{UNDEFINED} can be read in two ways. Under the
\emph{relaxed} reading (R1), the transition needs only an execution dependency on the write, since the old contents are
discarded. Under the \emph{strict} reading, it needs a memory dependency: the write must be made available before the
transition. The snapshots frozen for the tutorial study and the three engine studies applied R1 by default, while the
reference semantics of that period was strict, so their agreement with the reference holds only where R1 never
applies. The verified configuration is strict.

The specification does not allow R1. Its section on image layout transitions states that transitions ``may perform
read and write accesses on all memory bound to the image subresource range, so applications \emph{must} ensure that all
memory writes have been made available before a layout transition is executed'', with no exception for
\code{UNDEFINED}. \code{UNDEFINED} as an old layout only implies that ``the contents of the image subresource need not be
preserved'', a statement about the contents after the transition, not about ordering the transition's own writes.
Queue-family ownership transfers carry the same availability rule, and the Vulkan synchronization examples state that
a write-after-write through a transition from \code{UNDEFINED} still needs a memory dependency. The bad schedule R1
hides is a write-after-write on the image's memory: the earlier write, not made available, can land after the
transition has rewritten the subresource for its new layout.

\subsection{What the checker does}
\begin{itemize}
\item \textbf{Default (verified configuration): spec-strict.} A write before any layout transition must be made
  available by the transition's own barrier (first access scope: \code{srcStageMask} $\times$ \code{srcAccessMask} on
  that resource) or by a dependency chain whose first link does so.
\item \textbf{Opt-in heuristic.} An environment switch restores R1 for every harness; the report's metadata then
  records the heuristic and that the configuration is not verified.
\item \textbf{Diagnostics} name the failed scope (``is not made available before the layout transition \dots{} (from
  \code{UNDEFINED})'') and add, for \code{UNDEFINED}: ``the old contents may be discarded, but the earlier write must
  still be made available, or it can land after the transition.''
\item \textbf{Grading.} Report grouping grades a group a \emph{technicality} when every hazard in it concerns contents
  that are discarded anyway (the later access is a transition from \code{UNDEFINED}, or a load/store operation is
  \code{DONT\_CARE}). They are reported, but kept apart from errors.
\item \textbf{Final-layout transitions.} The trace translation starts a render pass's automatic transition to its
  final layout from the attachment's current layout. Hazard detection in the verified configuration does not depend
  on the old layout.
\end{itemize}
Nine unit tests cover the behaviour: the minimal case (write, barrier with \code{srcAccessMask = 0}, transition from
\code{UNDEFINED}) is reported by default and accepted under the heuristic; the default is recorded in the metadata; a
full memory dependency is clean; no execution dependency is a hazard in both modes; a defined old layout needs
availability in both modes; the access after the transition is still checked; a semaphore can carry the memory
dependency; and the reference defaults to strict.

\subsection{Effect on the earlier corpora}
Table~\ref{tab:s-undef} isolates the effect of strict transitions: the verified configuration against the heuristic
relaxation, both with conservative capture and both with the first snapshot of the verified configuration (post-hoc).
The difference is the same under observed capture on every corpus. That snapshot started final-layout transitions
from \code{UNDEFINED}; this does not change the strict column's hazard sets, but under the heuristic it made every
final-layout transition a discard, which the heuristic column includes.

\begin{table*}[t]
\caption{Effect of spec-strict transitions from \code{UNDEFINED}, first snapshot of the verified configuration\posthoc.
Engines: raw timing / report groups / technicality groups.}
\label{tab:s-undef}
\small
\begin{tabular}{@{}p{0.27\textwidth}p{0.14\textwidth}p{0.14\textwidth}p{0.37\textwidth}@{}}
\toprule
Corpus & strict (verified) & heuristic (R1) & effect of strict transitions \\
\midrule
\bench 114 (det. / unclean / exp. / missed) & 78 / 4 / 2 / 30 & 78 / 4 / 2 / 30 & none on status or code sets; raw counts differ in 3 buggy variants \\
Whole-frame, hazard / all-key agreement & 97.9\% / 95.4\% & 97.9\% / 95.4\% & 0 programs change \\
RT / mesh, 15 modes & 15 / 15 & 15 / 15 & 0 codes change; texts name the scope \\
Stored root causes (28 samples) & 24 IDs + T17 & 24 IDs + T17 & 0 lost; 356 diagnostics at new sites in 14 samples \\
\quad report groups error / technicality & 113 / 7 & 125 / 5 & merges in 7 samples \\
Tutorials (raw / groups) & 9,606 / 58 & 9,606 / 58 & none \\
SDL3 + bgfx & 221 / 21 / 0 & 221 / 21 / 0 & none \\
Granite, 4 & 589 / 4 / 8 & 519 / 6 / 6 & +70 raw, 16--19 per program, technicality \\
VSG, 8 & 471 / 9 / 1 & 431 / 9 / 0 & +40 raw, +1 technicality group (\code{vsg\_depthpeel}) \\
Diligent & 136 / 14 / 2 & 136 / 14 / 2 & none \\
vkQuake, 3 timedemos & 81,161 / 27 / 0 & 81,161 / 27 / 0 & none \\
\bottomrule
\end{tabular}
\end{table*}

\subsection{Triage of the diagnostics strict mode adds}
All added diagnostics are \code{timing.barrier-scope} of the form ``the earlier write (S, A) is not made available
before the layout transition \dots'', in three shapes.
\begin{itemize}
\item[(a)] \emph{A \code{DONT\_CARE} store or an attachment write, then a transition from \code{UNDEFINED}, with no
  \code{srcAccessMask} for the write.} Stored Sascha Willems samples (\code{deferred}, \code{deferredmultisampling},
  \code{deferredshadows}, \code{hdr}: 27 each; \code{ssao}, \code{subpasses}: 45 each; four others 18 each;
  \code{multiview} 9; \code{inputattachments} 27 of 47), Khronos \code{mobile\_nerf} (22), and all four Granite programs
  (16--19 each, on the depth prepass, after the previous frame's \code{DONT\_CARE} store). True by the specification,
  technicality in effect: the store operation is an attachment write, the transition needs it made available, and the
  effect is limited to contents that are then discarded.
\item[(b)] \emph{Ordering through a semaphore relay} (\code{synchronization\_2}, 8). Because a sync2 signal's stage
  mask limits every hop of a relay, the swapchain transition is the reported partner. The defect is KS-1, already in
  the list.
\item[(c)] \emph{Render-pass final-layout transitions} (\code{inputattachments} 20, \code{vsg\_depthpeel} 40). A
  \code{CLEAR} or \code{DONT\_CARE} write in the last subpass, then the automatic transition to the final layout, with a
  subpass-to-external dependency whose \code{srcAccessMask} lacks the depth write. True by the specification,
  technicality (\code{DONT\_CARE}). In this snapshot the messages of these 60 diagnostics say ``from
  \code{UNDEFINED}'' and the group's technicality reason names the transition rather than the \code{DONT\_CARE} store,
  because the translation then started final-layout transitions from \code{UNDEFINED}.
\end{itemize}

\paragraph{Group merges and site retention}
In four Sascha Willems samples, two error groups become one under strict transitions, because the synthesized repair
adds the write to the transition barrier's \code{srcAccessMask}, which also removes the downstream hazards. The
root-cause labelling rules of the sample study match message text, and the texts changed, so root causes are
compared by diagnostic site: every site that the engine-study-3 snapshot reports with a root-cause label is still
reported in every configuration, for each of the 24 labelled root-cause identifiers.

\subsection{VSG-5}
Synchronization validation reports 19 cross-frame write-after-write pairs on VSG-5 (the previous frame's depth store
against this frame's transition from \code{UNDEFINED}) that \sys does not report. They are not an effect of R1: strict
mode does not report them either, and \code{vsg\_rtt} gives the same 120 hazards strict or relaxed. \sys finds each
pair ordered by a chain through the window render pass's external depth dependency. This is the one disagreement with
validation in the accept direction that turns on a reading of the specification; the specification text decides it in \sys's favour, and the chain, the cited
clauses and a minimal reproduction are in \S\ref{sup:vsg5}. Every question we leave unresolved lies in the report
direction.

%% ===================================================================
\section{Checker Snapshots and Frozen Rounds}
\label{sup:history}

\subsection{What each study exposed}
Every study was scored by a checker snapshot frozen, by source hash, before its first run (Table~\ref{tab:s-versions}).
A checker defect found during a study was not fixed in that study's snapshot; it was fixed before a later snapshot,
with a unit test. Table~\ref{tab:s-exposed} lists what each study, and each piece of development work between the
studies, exposed in the checker, and the first snapshot that carries the fix. The release carries every fix in the
table and passes all 392 tests of its suite (paper, \S\ref{P-sec:eval-sem}).

\begin{table*}[t]
\caption{What each study exposed: checker defects and missing modelling found in triage or by tests, and the first
snapshot that carries the fix. Cause codes are those of \S F and \S I.}
\label{tab:s-exposed}
\scriptsize
\begin{tabular}{@{}p{0.2\textwidth}p{0.56\textwidth}p{0.18\textwidth}@{}}
\toprule
Study & defects found in triage & fixed before \\
\midrule
reference semantics (dev) & eight checker defects, each with a minimal frame on which the checker was silent or wrong, e.g.\ a layout transition not checked against later accesses; a semaphore wait chained through a barrier on another resource ignored, although execution dependencies are global; a wait with an empty stage mask taken to order something; \code{SHADER\_READ} taken to include \code{UNIFORM\_READ}; \code{ALL\_COMMANDS} in a later barrier taken to contain an earlier transition. Load/store operations became accesses, and external subpass dependencies are handled one at a time & tutorial-study snapshot \\
barrier mutants (dev; \S\ref{sec:eval-sem}) & two checker defects: access scopes stage-expanded like execution scopes; structs of one barrier call allowed to chain into each other & tutorial-study snapshot \\
sample collections and vuk (dev; \S F) & 11 systematic false-alarm causes (bindings declared but not used in an entry point, no array-layer tracking, page-granular host writes in sub-allocated memory, atomics treated as races, among others); ray tracing, mesh shading, events, rule R-LT and per-range ownership not yet modelled & during development (page granularity is again the default under conservative capture, \S C); modelling by blind round 6 \\
blind rounds 1--6 (fresh; \S\ref{sup:rounds}) & 11 timing cases missed entirely: synchronization not yet modelled (semaphore stage masks, ownership pairs, object lifetime, stencil-only aspects) and rules that did not fire; each miss became a general rule & the next round's snapshot \\
tutorials (fresh) & one false-alarm cause (FA-M: per-thread dispatches write disjoint ranges of one buffer; validation reports the same) & engine-study-2 snapshot \\
whole-frame differential test (fresh; \S\ref{sec:eval-wholeframe}) & six checker defects: a binary semaphore signalled again while a signal was pending not reported; a timeline signal whose value did not increase not reported; the sync2 stages \code{COPY}, \code{BLIT} and \code{CLEAR} merged into \code{TRANSFER}; a sync2 semaphore-signal stage mask ignored; command buffers, fences and acquire semaphores reused while pending not reported; grouping split one over-acquire site & engine-study-2 snapshot \\
engine study 1: vkQuake, Diligent (fresh) & false alarms from: shader index ranges not following push constants and dynamic offsets; depth accesses not following the depth-test state; inexact byte spans of truncated host-write records; late attribution of host writes; grouping split per recording command buffer. Missed (validation only): the first use of a never-presented swapchain image not ordered after its acquire (VD-2; rule \code{timing.acquire-first-use}) & engine-study-2 snapshot \\
engine study 2: SDL3, bgfx (fresh) & one checker defect behind all 13 false alarms (FA-WA): a semaphore wait with stage mask \code{ALL\_COMMANDS} not allowed to order a layout transition in the same batch, although the specification puts every command of the batch in the wait's second scope. Fixed, it removes exactly these 13 groups and two on one Khronos sample, and changes nothing on \bench, the tutorials, the whole-frame programs, vkQuake or Diligent. The same snapshot lets a \code{BOTTOM\_OF\_PIPE} semaphore signal cover an ownership release & engine-study-3 snapshot \\
engine study 3: Granite, VulkanSceneGraph (fresh) & one checker defect: push constants declared for \code{VK\_SHADER\_STAGE\_ALL} abort the check (CHK-PS; 3 programs unchecked). Missed (validation only): 16 write-after-write hazards on storage images bound through descriptor update templates, which are not modelled (a model limit, not fixed) & first snapshot of the verified configuration \\
formal model and re-examination of the earlier results (dev) & relaxed transitions from \code{UNDEFINED} by default (R1, \S D); observed host capture misses same-value stores (\S C); a sync2 signal stage mask did not limit every hop of a relay (third queue, host), in the checker and in the reference; a host-signalled timeline wait lost its stage mask; an atomic access could stand in for what it is ordered after; a shadow entry did not outlive a signal (found by a random frame); a binding charged at every stage its layout grants (\bench case 94); diagnostics did not name the failed scope test; elision compared diagnostic keys rather than event pairs & first snapshot of the verified configuration \\
post-hoc re-runs with that snapshot (\S D) & final-layout transitions of render passes started from \code{UNDEFINED} (in the verified configuration only messages and technicality reasons; under the heuristic, hidden hazards); repairs verified by diagnostic keys, so an existing witness could hide a new race & fresh-study-1 snapshot \\
fresh study 1 (fresh; \S I) & three checker defects: timeline counter query not treated as a host wait (FA-CQ); whole-image layouts as unbounded ranges (FA-WL); combined depth format \code{X8\_D24} treated as colour (FA-X8). Bring-up of the study's programs found one model-driver defect, fixed before the study's freeze; bring-up of study~2 needed \code{vkCmdClearAttachments} and \code{vkCmdCopyQueryPoolResults} & fresh-study-2 snapshot \\
fresh study 2 (fresh; \S I) & one precision cause: a descriptor-array read charged to every written element (GTK-DESC, \S B); grouping of the largest program exceeded the memory limit & the release \\
model audit after both fresh studies (dev) & subpass-dependency flags not read: every dependency given its global meaning (\S\ref{sup:fblocal}); two automatic render-pass transitions on one queue exempt from each other although the specification's ordering sentence does not cover them (\S\ref{sup:ltpairs}); in the epoch mode, the read key ignored the framebuffer position, so a stand-in could be ordered locally where the entry it replaced is not; host writes relying on the host-domain premise not flagged; guarantee eligibility not reported per run & the release \\
external review of the release (dev) & three defects where shader accesses become events: a store through an index taken from an atomic was exempt, although \code{atomicAdd(c, 0)} gives every caller one index; a missing shader module made storage accesses reads; an explicit relaxation was reported as the verified configuration. Fixed (slot exemption opt-in only, atomics exempt only at device scope, missing shader code a write and a model gap, one resolved configuration); findings unchanged on every corpus, and one archived figure trace kept without its shader module is now outside the domain & the release \\
second external review of the release (dev) & two shader-analysis defects that could accept a race: a helper function called with two buffers kept only one of them, so a write became a read; atomic-only tags compared block members, not buffer bytes, so an atomic at one descriptor offset hid a plain read of the same bytes at another. In the optimizer's split frame, the transition halves of two barriers could chain. Fixed (every possible descriptor kept, unattributed accesses a model gap, byte-range tags, transition halves chain only within one struct); findings unchanged on every corpus & the release \\
third external review of the release (dev) & the shader analysis read any instruction not in its list of writes as a read: \code{OpCopyMemorySized} (untyped pointers) and cooperative-matrix stores wrote without a write event, and an untyped-pointer access chain was not reported. Fixed: an inventory of memory effects (Table~\ref{tab:s-effects}); any other instruction on descriptor memory is a model gap. On every shader of every corpus (2,129 entry points) the analysis is unchanged & the release \\
fourth external review of the release (dev) & the byte span of an access through \code{OpPtrAccessChain} took the chain's Element operand for the block member, so a plain read could be placed on other bytes than it touches and escape an atomic it overlaps (a module that passes \code{spirv-val}); the size of a row-major non-square matrix could be underestimated. Fixed (the Element is skipped and must be 0, otherwise the whole range; matrix size bounded for both layouts); unchanged on every corpus shader & the release \\
fifth external review of the release (dev) & an atomic whose memory scope is a specialization constant, which the KHR and NV cooperative-matrix extensions allow, was classified by the constant's default, so a pipeline that specializes Device to Workgroup kept the device-scope exemption; a specialized pointer-chain Element or member index was likewise read by its default. Fixed (only \code{OpConstant} values are read; otherwise an ordinary access or the whole range), with a regression whose NV form passes \code{spirv-val}; the access classification is unchanged on all 259 distinct shader modules of the corpora & the release \\
\bottomrule
\end{tabular}
\end{table*}

\subsection{The frozen blind rounds}
\label{sup:rounds}
\paragraph{Protocol}
Before each round the checker was frozen and the outcome of every case was predicted and timestamped before the run.
Round 1 was written after the checker was frozen. Round 2 was written by a separate party that never ran the checker.
Rounds 3 and 4 were written without reading any checker code; the case author knew only the list of diagnostic codes.
Rounds 5 and 6 each used 20 bugs collected fresh from public trackers, each with a merged upstream fix, built blind
without reading checker code (round 6 also without reading the paper); round 6 ran on the checker with ray tracing,
mesh shading and R-LT. After each round, misses became general rules and the round's cases became development data.

\paragraph{Metrics}
\emph{Strict paired-detection success}: the buggy variant reports every expected code and the fixed variant is clean.
\emph{Root-cause detection}: the buggy variant reports at least one timing diagnostic on the case's root-cause resource,
whatever its code. \emph{Fixed-variant diagnostics}: any diagnostic on the fixed variant. The root-cause keys were fixed
from each case's expected finding and debug names before any post-hoc result was seen. Intervals are Wilson 95\%.

\begin{table*}[t]
\caption{Blind rounds, as run (fresh): each round scored by the checker frozen for it.}
\label{tab:s-rounds}
\small
\begin{tabular}{@{}lrrrrrrr@{}}
\toprule
Round (cases) & cases & timing & strict (timing) & root cause (timing) & fixed, any diagnostic & fixed, timing & timing raw on fixed \\
\midrule
1 (17--29) & 13 & 7 & 3 & 4 & 1 & 1 & 3 \\
2 (30--44) & 15 & 7 & 4 & 4 & 2 & 0 & 0 \\
3 (45--59) & 15 & 6 & 3 & 3 & 0 & 0 & 0 \\
4 (60--74) & 15 & 5 & 3 & 4 & 8 & 5 & 12 \\
5 (75--94) & 20 & 6 & 5 & 5 & 2 & 1 & 7 \\
6 (95--114) & 20 & 4 & 3 & 4 & 1 & 0 & 0 \\
\midrule
total & 98 & 35 & 21 (60.0\%, 43.6--74.4) & 24 (68.6\%, 52.0--81.4) & 14 (14.3\%, 8.7--22.6) & 7 (7.1\%, 3.5--14.0) & 22 \\
\bottomrule
\end{tabular}
\end{table*}

Root cause found but not strict: 28 and 63 (detected, fixed variant not clean) and 98 (the host race on the decoder
buffer is reported, under another code). Timing cases missed entirely: 17, 19, 20 (known misses) and 32, 36, 38, 45, 47,
50, 61, 78. Fixed variants with a timing diagnostic: 28 and 63 (timing cases); 65, 66, 68 and 69 (round-4 meaning cases;
both variants report the same host-write race in the shared harness); and 94 (a structure case; a binding charged at a
stage whose entry point does not use it). Predictions held in 13/13, 13/15, 14/15, 13/15, 17/20 and 14/20 cases. Round
5, on 20 bugs from 16 projects (13 new to the benchmark), found 5 of 6 timing bugs; the miss needs host-side availability
after a fence wait. In all categories, strict detection was 29 of 98 (30\%, 21--39\%).

\paragraph{Post-hoc re-run in the verified configuration}
All 114 cases were built, traced and checked with the first snapshot of the verified configuration; the other
configurations, and the engine-study-3 snapshot, re-check the same traces, so the checker is the only variable
(Table~\ref{tab:s-roundsph}). The timing columns are the same in every row: the same 29 strict detections and 33 root
causes. The step from 21 and 24 as run to 29 and 33 is development between rounds, made after each round was seen, not
an effect of the configuration. Remaining strict misses: 47, 50, 61, 62, 78, 98; root-cause misses: 50, 78. Fixed
variants: 71 and 73 draw the same convention warning everywhere; 61 is a \code{signal-stage-scope} on the present
semaphore outside the case's fix; 62 is the first-acquire technicality; 94 (engine-study-3 snapshot only) was the
stage-attribution false alarm; 65, 66, 68 and 69 (conservative capture) are the true harness race. The release gives
the same statuses and codes on all 114 cases.

\begin{table}[htbp]
\caption{Blind timing cases and fixed variants, re-run on traces recorded for the first snapshot of the verified
configuration (post-hoc), per configuration; the first row re-checks the same traces with the engine-study-3
snapshot.}
\label{tab:s-roundsph}
\scriptsize
\resizebox{\columnwidth}{!}{%
\begin{tabular}{@{}lrrrrr@{}}
\toprule
Configuration & strict & root cause & fixed, any & fixed, timing & timing raw \\
\midrule
engine-study-3 snapshot & 29 & 33 & 5 & 3 (61, 62, 94) & 18 \\
\textbf{verified} & 29 & 33 & 8 & 6 (61, 62, 65, 66, 68, 69) & 24 \\
observed capture & 29 & 33 & 4 & 2 (61, 62) & 16 \\
heuristic R1 & 29 & 33 & 8 & 6 & 24 \\
heuristic R1, observed capture & 29 & 33 & 4 & 2 & 16 \\
\bottomrule
\end{tabular}}
\end{table}

%% ===================================================================
\section{Per-Study and Per-Engine Triage}
\label{sup:triage}

\paragraph{Common protocol}
For each fresh study, the checker and a triage protocol were frozen before the first scored run. Every timing report
group was triaged against the specification and the source as TRUE (a real bug), TECH (a real violation without
practical effect, e.g.\ on contents discarded anyway, or a universally tolerated pattern), FA (false alarm), DRV (an
artefact of the model driver; excluded from denominators) or UNRESOLVED. \code{timing.unmodelled-command} warnings are
listed but not scored. Synchronization validation (all synchronization checks on, lavapipe) on the
same frames is supporting evidence only.

\subsection{Tutorials (fresh)}
All 36 chapter programs of Khronos' Vulkan-Tutorial and the 7 chapters of vulkan-guide
built and ran unmodified for 10 frames; 29 record frame work and were checked. \sys reported 25 report
groups in 8 programs (Table~\ref{tab:s-tut}).

\begin{table*}[t]
\caption{Tutorial study, snapshot frozen for it (fresh).}
\label{tab:s-tut}
\small
\begin{tabular}{@{}llrrp{0.5\textwidth}@{}}
\toprule
Cause & class, code & groups & raw & evidence \\
\midrule
VG-1 & TRUE, stale-barrier & 2 & 11 & chapter 5 transitions the draw image with a stale \code{oldLayout} (fixed upstream in chapters 3--4 only) \\
VG-2 & TRUE, missing-barrier & 3 & 3,135 & chapter 6: compute writes the draw image, rendering then loads and writes it, no barrier \\
VG-3 & TRUE, wrong-layout & 6 & 62 & chapter 6 renders the UI into the swapchain image declared \code{GENERAL} \\
TECH-S & TECH, present-semaphore-reuse & 13 & 63 & render semaphore per frame in flight, re-signalled after the frame fence only \\
FA-M & FA, missing-barrier & 1 & 70 & per-thread dispatches write disjoint ranges of one buffer; validation reports the same 70 \\
\bottomrule
\end{tabular}
\end{table*}
Result: 11 / 13 / 1 of 25 groups (44\% true, 96\% actionable); by raw diagnostic 96.0\% true; by cause 3 / 2 / 1.
Validation produced no device-level message on any vulkan-guide program for a reason not found (a 60-line reproduction
of VG-1 is reported by the same validator build), so VG-1 to VG-3 rest on source, trace and specification. Under the
engine-study-2 and engine-study-3 snapshots (dev), grouping gives 6 / 13 / 0 of 19 with no finding lost.

\subsection{vkQuake and Diligent (fresh)}
vkQuake (three timedemos, 107--182 presents) and Diligent's glTF viewer (20 frames).
\begin{itemize}
\item \textbf{TRUE (VQ-1 + VQ-2)}, one repair group per demo: the indirect-argument and index buffers exist once and the
  next frame's compute rewrites them while the previous frame's indirect draws may read them (VQ-1; validation: 39 / 39
  / 19 write-after-read); and the scene render pass has no outgoing dependency, so final transitions of the shared depth
  and transparency attachments are unordered with the next frame (VQ-2; validation reports the mechanism on another
  attachment). One synthesized repair covers both, so they share a group.
\item \textbf{TRUE (VD-1)}, five groups: Diligent's state tracker does not count depth writes as writes, so
  consecutive dynamic-rendering instances get no depth barrier (validation: 9 read-after-write, 5 write-after-write).
\item \textbf{FA}: host writes into per-frame halves of ring buffers charged to the whole buffer and split into one group
  per recording command buffer (FA-R, the bulk); disjoint storage-image regions of one dispatch sequence (FA-M, shared
  with validation); a depth write charged with the depth test disabled (FA-DT); a page fallback on a truncated host
  record (FA-PG); and three non-timing causes.
\item \textbf{DRV}: one cause (late attribution of host writes), excluded.
\item \textbf{Validation only}: VD-2, the first transition of a never-presented swapchain image not ordered after the
  acquire; the first-acquire rule of the engine-study-2 snapshot closes it.
\end{itemize}
Result (timing): 8 / 0 / 44 of 52 groups (15.4\%); 97.9\% of raw diagnostics true, because one true group carries
almost all of them; by cause 3 / 0 / 7. Under the engine-study-2 and engine-study-3 snapshots (dev): 7 / 2 / 11 of 20.

\paragraph{Verified configuration (post-hoc)}
On newly recorded traces, the engine-study-3 snapshot and the first snapshot of the verified configuration with
observed capture give 7 / 2 / 11 of 20 (vkQuake 3 / 0 / 11 of 14, Diligent 4 / 2 / 0 of 6). The verified configuration
gives 7 / 2 / 34 of 43 (vkQuake 3 / 0 / 24 of 27, Diligent 4 / 2 / 10 of 16); the 23 new
groups are conservative-capture false alarms classified in Table~\ref{tab:s-capcost}. Every vkQuake timedemo also
carries one P3 warning (worker-thread host writes).

\subsection{SDL3 and bgfx (fresh)}
SDL3 (two Vulkan back ends, 6 programs) and bgfx (6 programs), 20 presents each.
\begin{itemize}
\item \textbf{TRUE (SDL-1)}, 2 groups (two directions of one race, 166 raw): SDL's Vulkan renderer keeps its vertex
  buffers globally rather than per frame, restarts at index 0 every frame and copies into the mapped buffer with no
  wait; the only wait is on the fence of the frame \code{swapchainImageCount} frames back, while later frames that
  draw from the same buffer are pending. Validation does not see host writes.
\item \textbf{FA (FA-WA)}, 13 groups: one per SDL program with the defect and two per bgfx program
  (\code{wait-stage-scope} and \code{acquire-first-use}). A semaphore wait with stage mask \code{ALL\_COMMANDS} was not
  allowed to order a layout transition in the same batch. Validation is silent, matching the specification.
\end{itemize}
Result: 2 / 0 / 13 of 15 (13.3\%); raw 57.2\% true; by cause 1 / 0 / 1. Under the engine-study-3 snapshot (dev):
2 / 0 / 0 of 2; in every configuration of the first snapshot of the verified configuration (post-hoc): the same. Structure and meaning (outside the paper's claims): 17 groups, all false
alarms, four causes.

\subsection{Granite and VulkanSceneGraph (fresh)}
Granite (4 programs) and VulkanSceneGraph with vsgExamples (8 programs); 3 programs
aborted under the snapshot frozen for the study (CHK-PS), so 9 were checked.
\begin{itemize}
\item \textbf{VulkanSceneGraph, 8 groups, all TRUE, 5 causes.} VSG-1 (library): the semaphore wrapper's default wait
  stage is \code{BOTTOM\_OF\_PIPE}, so the consumer-completed semaphores order nothing (validation: 18
  write-after-read per program). VSG-2 (library): two consecutive transfer submits write the same destination buffers
  with nothing ordering them (validation silent; on lavapipe's single queue the submissions differ). VSG-3 (example):
  a compute write to a vertex buffer and the vertex-input read are separated only by a colour-only subpass dependency
  (validation: 19 read-after-write). VSG-4 (example): a pre-copy barrier starts at \code{TOP\_OF\_PIPE}, an empty first
  scope (validation silent on this pair). VSG-5 (example): a depth clear is not ordered after the render pass's own depth
  layout transition (validation: 19 + 19 write-after-write; validation's further 19 cross-frame pairs are discussed in
  \S\ref{sup:vsg5}).
\item \textbf{Granite, 12 groups, all TECH, one cause.} The deferred-lighting pass takes the G-buffer as read-only input
  attachments and still declares \code{storeOp STORE}; the specification makes the store an attachment write, so the
  write-after-write needs a memory dependency, but nothing writes those attachments during the pass, so the store
  writes back the bytes already there.
\end{itemize}
Result: 8 / 12 / 0 of 20 (40.0\%; 100\% actionable); by cause 5 / 1 / 0; raw 334 / 516 / 0 (39.3\%). Validation's 16
write-after-write hazards in Granite's ambient-occlusion chain, on storage images bound through descriptor update
templates (not modelled), are a recall gap of \sys. The three windowed Granite programs crash on lavapipe, so only the
headless one has a validation baseline.

\paragraph{Verified configuration (post-hoc)}
With CHK-PS fixed (first snapshot of the verified configuration), all 12 programs run. Granite: 4 programs, each 1
error and 2 technicality groups, all TECH (the same root plus the strict-\code{UNDEFINED} pairs on the depth prepass).
VulkanSceneGraph, the five programs checked before: 8 groups, hazard pair sets identical to the study's snapshot
(epoch 151 = 151, pairwise 224 = 224), all TRUE. The three that the study's snapshot could not check:
\code{vsg\_gltf} and \code{vsg\_shadow} report nothing; \code{vsg\_depthpeel} reports one technicality group (40 raw,
render-pass final-layout transitions, \S D) and one error group we leave unresolved (60 raw). Its window render pass
ends with the implicit dependency, whose second scope is empty, and a second render pass in the same command buffer then
loads and writes the swapchain image. \sys reads the final transition to the present layout as unordered with those
accesses; validation accepts them (Table~\ref{tab:s-vsg5}: also when the final transition is done by a pipeline barrier
with \code{dstStageMask = BOTTOM\_OF\_PIPE}), so validation lets a later dependency's first scope include an earlier
transition whose own dependency has an empty second scope, which the cited clauses do not say. With an explicit exit
dependency at \code{COLOR\_ATTACHMENT\_OUTPUT}, both tools are clean. The question needs a ruling; it is in the report
direction, and we do not count the group as a bug. Result: 8 / 13 / 0 / 1 of 22.

\subsubsection{VSG-5: the disagreement in the accept direction}
\label{sup:vsg5}
Synchronization validation reports 19 cross-frame write-after-write pairs on VSG-5 (\code{vsgrendertotexture}): the
previous frame's depth store against this frame's transition of the depth image from \code{UNDEFINED}. \sys finds
each pair ordered, in the strict and in the relaxed reading of \S D. Of the disagreements with validation that turn on
a reading of the specification, this is the one in the accept direction (the other pairs that validation alone reports
are unmodelled API features, driver differences, or findings both tools report under different keys; \S F and \S H): here \sys's verdict is the more permissive one,
and a clean verdict could rest on it.

\emph{The chain.} (1) The library's window render pass, recorded after the example's offscreen pass in the same command
buffer, has an \code{EXTERNAL}$\to$0 depth dependency with source stages
\code{EARLY\_FRAGMENT\_TESTS|LATE\_FRAGMENT\_TESTS} and source access \code{DEPTH\_STENCIL\_ATTACHMENT\_WRITE}, with no
resource restriction; the store of frame $k$ is in its first access scope. (2) The next frame's offscreen pass has an
\code{EXTERNAL}$\to$0 dependency with \code{srcStageMask FRAGMENT\_SHADER}, whose first synchronization scope includes
the logically earlier \code{EARLY\_FRAGMENT\_TESTS}; this dependency carries the transition. (3) The chain condition
holds through \code{EARLY\_FRAGMENT\_TESTS}, so the store happens before the transition and is available to it.

\emph{The specification} (Vulkan-Docs v1.4.365) decides the case. In the render pass chapter, external
subpass dependencies: with \code{srcSubpass = VK\_SUBPASS\_EXTERNAL}, ``the first synchronization scope includes
commands that occur earlier in submission order than the \code{vkCmdBeginRenderPass}'', limited only by
\code{srcStageMask} (renderpass.adoc 3693--3706); the first access scope ``is limited to accesses in the pipeline stages
determined by the source stage mask'' and to the access types in the source access mask (3723--3728), and nothing
limits it to the render pass's own attachments; the transition away from \code{initialLayout} happens after the
availability operations of every \code{EXTERNAL}$\to$\code{dstSubpass} dependency (3975--3979). The chain condition of
synchronization.adoc (94--96: the second synchronization scope of one dependency intersects the first synchronization
scope of the next) then holds through \code{EARLY\_FRAGMENT\_TESTS}. With these clauses the chain exists and the 19 pairs
are ordered.

\emph{Minimal reproduction.} One headless program on one queue, render passes with load and store operations only (no
draws, no shaders), with every frame recorded into one command buffer, so that consecutive frames are ordered only by
what they record (Table~\ref{tab:s-vsg5}). Validation's verdict on the cross-frame pair changes from 2 reports to 0
when the window pass's dependency is replaced by a pipeline barrier with identical stage and access masks, while the
specification gives the two the same first synchronization and access scopes. This suggests that validation applies an
external subpass dependency only to the accesses of the pass's own attachments. Both tools agree on VSG-5's main finding
(the depth clear is not ordered after the depth layout transition) in every variant.

\begin{table}[htbp]
\caption{Minimal reproductions of the two disagreements with validation (validation on lavapipe; \sys:
fresh-study-1 snapshot). VSG-5: load-operation pairs + store-operation pairs, and cross-frame pairs.
\code{vsg\_depthpeel}: pairs of pass~1's final transition with pass~2's load and store.}
\label{tab:s-vsg5}
\scriptsize
\resizebox{\columnwidth}{!}{%
\begin{tabular}{@{}lll@{}}
\toprule
variant & validation & \sys \\
\midrule
\multicolumn{3}{@{}l}{\emph{VSG-5: offscreen depth pass A, then the window pass B, per frame}} \\
as recorded & 3 + \textbf{2 cross-frame} & 3 + 3, \textbf{0 cross-frame} \\
pass B removed & 3 + 2 cross-frame & 3 + 3, cross-frame reported \\
pass B replaced by a barrier with the same masks & 3, \textbf{0 cross-frame} & 3 + 3, 0 cross-frame \\
\midrule
\multicolumn{3}{@{}l}{\emph{\code{vsg\_depthpeel}: pass 1 with the implicit exit dependency, then pass 2}} \\
as recorded & \textbf{0} & \textbf{2} \\
explicit \code{COLOR\_ATTACHMENT\_OUTPUT} exit dependency & 0 & 0 \\
final transition by a barrier to \code{BOTTOM\_OF\_PIPE} & \textbf{0} & \textbf{2} \\
\bottomrule
\end{tabular}}
\end{table}

\emph{Ablation on the recorded trace} (first snapshot of the verified configuration). Unablated, \code{vsg\_rtt} gives
120 diagnostics in both readings of \S D. With the 20 window-pass dependencies removed from the model, the strict
reading gives 716 diagnostics (317 with a store-operation partner) and the relaxed one 678 (298), so the strict reading
reports 19 store-operation to transition pairs that the relaxed one does not: without the window pass's dependency,
\sys reports the pairs that validation reports.

\emph{Consequence.} The cited text decides VSG-5 in \sys's favour, so we count validation's 19 cross-frame pairs as
probable validation false positives; VSG-5's main finding (20 load-operation versus transition pairs) is reported by
both tools. No accepted ordering in our results rests on an unresolved reading of the specification: every question we
leave unresolved, \code{vsg\_depthpeel} (above and Table~\ref{tab:s-vsg5}) and the pairs of automatic render-pass
transitions in GTK and Qt (\S I), lies in the report direction. There \sys reports the pair and validation accepts it,
so the question can only add reports and no clean verdict depends on it.

\subsection{Development data: two sample collections and a render graph}
SaschaWillems/\-Vulkan (all 99 samples) and Khronos\-Group/\-Vulkan-\-Samples (the 49 samples
that record work under the study driver), 10 frames each, plus 27 ray-tracing, mesh-shading and compute samples and the
78 tests of the vuk render-graph library. Every report was triaged. The checker was then fixed using
these reports, so they are development data (Table~\ref{tab:s-real}).

\begin{table}[htbp]
\caption{Triage of reports on unmodified samples (study snapshot; dev).}
\label{tab:s-real}
\small
\begin{tabular}{@{}lrrrr@{}}
\toprule
& \multicolumn{2}{c}{Sascha Willems (99)} & \multicolumn{2}{c}{Khronos (49)} \\
\cmidrule(lr){2-3}\cmidrule(l){4-5}
Class & causes & groups & causes & groups \\
\midrule
true bug & 17 & 63 & 5 & 20 \\
technicality & 3 & 31 & 3 & 87 \\
false alarm & 11 & 53 & 6 & 164 \\
driver / out of scope & 2 & 7 & -- & -- \\
\bottomrule
\end{tabular}
\end{table}

The 22 root causes judged real at the study snapshot include six host-write races (T1--T6), compute and graphics
submissions sharing a storage resource with no semaphore (T7), swapchain transitions from \code{TOP\_OF\_PIPE} not ordered
after the acquire semaphore (T9a), offscreen passes whose outgoing dependency ends at \code{BOTTOM\_OF\_PIPE} (T14, six
samples), a host struct 24 bytes long against a std430 stride of 32 (T17), and a semaphore signal with
\code{stageMask = 0} (KS-1). Four more were found afterwards, for 26: SW-13 (all nine Sascha Willems ray-tracing samples)
and KS-7 (three Khronos ones), where the ray-traced storage image is transitioned with no source access; KS-6, a
fragment buffer allocated with 12-byte elements against a stride of 16; and VUK-1, a bug in vuk's own barrier
derivation. Validation reports 10 of the Sascha Willems and 4 of the Khronos root causes at the study snapshot, and
confirms SW-13, KS-7 and, on a variant that uploads with a GPU copy, VUK-1; it cannot see the host races. The false
alarms came from 11 systematic causes (bindings declared but not used in an entry point, no array-layer tracking,
page-granular host writes in sub-allocated memory, atomics treated as races, among others), all since fixed with unit
tests. On the Sascha Willems T2--T5 races, see \S\ref{sup:dossiers}.

%% ===================================================================
\section{Bug Dossier Summary}
\label{sup:dossiers}

We judge 38 bugs to be real: 26 in the sample collections and vuk, 3 in the tutorials and 9 in the engines. This triage
is ours. For the 11 strongest we prepared an evidence dossier, the \emph{11 historical dossiers}: the affected
revision, the code, the diagnostic of the engine-study-3 snapshot verbatim, validation's messages on the same frames
where validation can see the bug, a specification argument with a concrete bad schedule, a patch, and logs before and
after under both tools. They are distinct from the 11 source-patch cases of the fresh studies (\S\ref{sup:repairs}),
which cover nine causes; SDL-1 is in both. All dossier runs use the engine-study-3 snapshot and
its model driver; the validation baseline is lavapipe with the validation commit of Table~\ref{tab:s-software}. These are \emph{source repairs} (the third kind of repair evidence in the paper's \S\ref{P-sec:eval-bugs}). A source repair counts as \emph{re-checked} when,
after rebuilding and re-running, \sys's count on the bug is 0, validation's count is 0 where validation can see the bug,
and no new diagnostic kind appears. The affected code was unchanged at each project's head when its dossier was
written. None has been confirmed by a maintainer yet.

\begin{table*}[t]
\caption{The 11 historical dossiers. Syncval: validation reports the same resource and command pair (n/a: a host-write
race, which validation does not track). Repair: \sys count (engine-study-3 snapshot) before $\to$ after; validation in
parentheses.}
\label{tab:s-dossiers}
\small
\begin{tabular}{@{}lllp{0.33\textwidth}lr@{}}
\toprule
Bug & Project & Revision & Class & Syncval & Repair \\
\midrule
KS-1 & Vulkan-Samples & \code{177edebf} & semaphore signal with \code{stageMask = 0} (empty first scope) & yes (11) & 99 $\to$ 0 (11 $\to$ 0) \\
KS-3 & Vulkan-Samples & \code{177edebf} & no barrier between two dynamic-rendering instances on the swapchain image & yes (8) & 40 $\to$ 0 (8 $\to$ 0) \\
SW-T1 & Sascha Willems Vulkan & \code{41a4410} & host write to a uniform buffer still read by a frame in flight & n/a & 500 $\to$ 0 \\
SW-T6 & Sascha Willems Vulkan & \code{41a4410} & host write to a uniform buffer right after the compute submit that reads it & n/a & 2 $\to$ 0 \\
SDL-1 & SDL3 & \code{e4af6f8} & host write to vertex buffers shared by frames in flight & n/a & 166 $\to$ 0 \\
VUK-1 & vuk & \code{61abde9} & render-graph barrier derivation: storage-read access for a uniform buffer & yes (1) & 1 $\to$ 0 (1 $\to$ 0) \\
VSG-1 & VulkanSceneGraph & \code{2feabc6} & semaphore waited at \code{BOTTOM\_OF\_PIPE} (empty second scope) & yes (18) & 19 $\to$ 0 (18 $\to$ 0) \\
VSG-3 & vsgExamples & \code{46a82943} & no barrier between compute write and vertex-input read & yes (19) & 58 $\to$ 0 (19 $\to$ 0) \\
VSG-5 & vsgExamples & \code{46a82943} & subpass dependency without depth stages: depth clear vs.\ layout transition & yes (20) & 120 $\to$ 0 (39 $\to$ 0) \\
VD-1 & DiligentCore & \code{074060c} & state tracker does not treat depth writes as writes & yes (9 + 5) & 114 $\to$ 0 (14 $\to$ 0) \\
VQ-1 & vkQuake & \code{749b4fd} & GPU-written indirect buffers shared by frames in flight, no barrier & yes (19) & 302 $\to$ 0 (19 $\to$ 0) \\
\bottomrule
\end{tabular}
\end{table*}

\paragraph{Notes on the dossiers}
\begin{enumerate}
\item \emph{Sascha Willems T2--T5} (host-write races in \code{particlesystem}, \code{rayquery}, \code{gltfskinning} and
  \code{computecloth}) are real in the source, but the engine-study-2 and engine-study-3 snapshots did not detect them
  on newly recorded traces: the racing stores write the bytes already there, and observed capture charges only changed
  bytes. They were not given dossiers. The traces stored at the sample study and conservative capture (post-hoc) report
  them (T2 8, T3 4, T4 9, T5 576 raw). T6 has only 2 diagnostics under the engine-study-3 snapshot for the same reason,
  and T1's view uniform buffer the same blind spot.
\item \emph{VSG-3 and VSG-5 are in vsgExamples}, not in the VulkanSceneGraph library: VSG-3 is a missing barrier in the
  compute-vertex example and VSG-5 a hand-written dependency in the render-to-texture example. Per project:
  VulkanSceneGraph 2 (VSG-1, VSG-2), vsgExamples 3 (VSG-3, VSG-4, VSG-5).
\item \emph{VSG-5's 19 cross-frame pairs}, reported by validation but not by \sys, are not an effect of R1 (strict
  transitions do not report them either); the specification decides them in \sys's favour (\S\ref{sup:vsg5}).
\item \emph{KS-1} also contains a second, independent true bug (a compute-to-compute write-after-write behind a barrier
  whose destination access is read-only; 11 raw), corroborated by validation's one write-after-write; it is an optional
  second patch in the dossier.
\item \emph{VQ-1 shares a report group with VQ-2}, because one synthesized repair covers both; the dossier counts raw
  diagnostics per resource (VQ-1 302, VQ-2 about 22,500).
\item \emph{KS-3's raw count} depends on the snapshot (18 raw at the sample study, 40 under the engine-study-3
  snapshot), because the latter also lists the write-after-write pairs into the UI pass; same root cause and fix.
\end{enumerate}
Not dossiered: VD-2, VSG-2, VSG-4 and VQ-2 (VQ-2 is real, but its validation match is the same mechanism on another
attachment and its fix is larger), T2--T5, and the remaining sample root causes.

%% ===================================================================
\section{Full Results Tables}
\label{sec:eval}

\subsection{\bench}
\bench has 166 bugs from 105 projects: the first collection of 126 bugs from 79 projects (48 timing, 27 structure, 51
meaning; Table~\ref{tab:s-taxonomy}) and 40 bugs collected fresh for rounds 5 and 6 (10 timing, 9 structure, 21
meaning), each with a merged upstream fix. 64 first-collection rows link a fix and 48 are both fixed and linked. About
60 of the 79 projects are small; the 29 well-known projects contribute 74 bugs. Table~\ref{tab:s-benchver} gives the
strict status of all 114 reproductions per checker snapshot and configuration.

\begin{table}[htbp]
\caption{\bench taxonomy (first collection, 126 bugs). ``Fixed'': fixed with a linked fix.}
\label{tab:s-taxonomy}
\scriptsize
\begin{tabular}{@{}llrr@{}}
\toprule
Category & Subcategory group & Bugs & Fixed \\
\midrule
\multirow{8}{*}{Timing (48)} & swapchain / semaphore misuse & 15 & 5 \\
 & missing barrier / wrong stage or access & 10 & 2 \\
 & image layout transition & 8 & 1 \\
 & cross-queue / ownership & 5 & 1 \\
 & frames-in-flight reuse & 4 & 2 \\
 & render-pass / subpass dependency & 3 & 2 \\
 & resource state (D3D12) & 2 & 1 \\
 & barrier in loop & 1 & 0 \\
\midrule
\multirow{6}{*}{Structure (27)} & host/shader layout, padding & 9 & 5 \\
 & descriptor binding / set & 8 & 3 \\
 & vertex input / stage interface & 5 & 1 \\
 & push-constant range / stage & 3 & 1 \\
 & bindless / descriptor heap & 1 & 1 \\
 & framebuffer / attachment extent & 1 & 0 \\
\midrule
\multirow{9}{*}{Meaning (51)} & color encoding / format & 26 & 10 \\
 & depth convention & 9 & 6 \\
 & handedness & 3 & 2 \\
 & Y-flip / clip origin & 3 & 0 \\
 & premultiplied alpha & 3 & 2 \\
 & coordinate frame & 2 & 1 \\
 & row- vs.\ column-major & 2 & 0 \\
 & value range & 2 & 1 \\
 & units & 1 & 1 \\
\bottomrule
\end{tabular}
\end{table}

\begin{table}[htbp]
\caption{Strict status of all 114 \bench reproductions per snapshot and configuration: detected / detected with unclean fixed variant /
expected miss / missed. All are development data for cases 1--16 and re-runs (not the blind per-round result) for
17--114.}
\label{tab:s-benchver}
\scriptsize
\resizebox{\columnwidth}{!}{%
\begin{tabular}{@{}llll@{}}
\toprule
Snapshot & status & provenance & mode \\
\midrule
tutorial-study snapshot & 81 / 1 / 2 / 30 & dev & R1, observed \\
engine-study-2 snapshot & 80 / 2 / 2 / 30 & dev & R1, observed \\
engine-study-3 snapshot & 80 / 2 / 2 / 30 & dev & R1, observed \\
\quad the same, newly recorded traces & 80 / 2 / 2 / 30 & post-hoc & R1, observed \\
first verified-configuration snapshot & 78 / 4 / 2 / 30 & post-hoc & \textbf{verified} (strict, conservative) \\
\quad observed capture & 81 / 1 / 2 / 30 & post-hoc & strict, observed \\
\quad heuristic R1 & 78 / 4 / 2 / 30 & post-hoc & R1, conservative \\
\quad heuristic R1, observed capture & 81 / 1 / 2 / 30 & post-hoc & R1, observed \\
\bottomrule
\end{tabular}}
\end{table}

The engine-study-3 snapshot on the newly recorded traces reproduces its stored statuses exactly, so the trace driver
changes no status. From that snapshot to the verified configuration on the same traces: 65, 66 and 68 become
fixed-not-clean and 69 gains a timing code on both variants (conservative capture); 94 becomes a clean detection and 42
loses a spurious timing code (entry-point stage attribution); 20 gains a \code{signal-stage-scope} diagnostic (signal
stage masks limit every hop of a relay). After round 6, under strict scoring, the tutorial-study snapshot detects 35 of
the 40 timing cases, and the engine-study-3 snapshot and every configuration of the first verified-configuration
snapshot 34, because the fixed variant of case 62 draws the first-acquire warning added in the engine-study-2 snapshot.
These are tuned numbers. The release gives the same statuses and codes as the verified
row.

\subsection{Semantics checks}
\label{sec:eval-sem}

\begin{table}[htbp]
\caption{Checker vs.\ executable reference semantics, early checker around the tutorial-study snapshot (dev), R1 in
both. ``Violations'': reference races the checker misses. The release against the reference is in the paper
(\S\ref{P-sec:eval-sem}).}
\label{tab:ref}
\small
\resizebox{\columnwidth}{!}{%
\begin{tabular}{@{}lrrrr@{}}
\toprule
Input & Frames & Ref.\ races & Checker & Violations \\
\midrule
random frames (9 option sets $\times$ 150 seeds) & 1,350+ & -- & -- & 0 \\
difftest mutants & 753 & 6,279 & 6,279 & 0 \\
bench traces (bug + fixed) & 148 & 709 & 709 & 0 \\
elide traces, base and with plans & 6 & 0 & 0 & 0 \\
\bottomrule
\end{tabular}}
\end{table}

\paragraph{Reference semantics}
On every input there is no race that the reference finds and the checker misses; the checker is stricter only on the
option set that exercises partial region coverage (the region cut of Proposition~\ref{P-thm:complete}). The reference
shares the checker's reading of the specification. At the time of this table it also shared two gaps, fixed since in
both: a sync2 signal stage mask did not limit every hop of a relay, and a host-signalled timeline wait lost its stage
mask (\S E). It gained the mesh, ray-tracing,
acceleration-structure and event pipelines; random event frames agree exactly (150 seeds in the test suite, 3,000 once),
and disabling the event first-scope cut makes 158 of them fail.

\paragraph{Barrier mutants}
750 random barrier plans (250 per workload: \emph{particles}, \emph{deferred}, \emph{postchain}) remove or narrow
barrier structs within the recorded masks. On the validation commit pinned for every later study (full
synchronization settings, lavapipe) and the fresh-study-1 snapshot (post-hoc), mutant-level agreement is 95.9\%: 86
both safe, 633 both hazard, 31 checker only and 0 validation only; at the resource level, 2,616 both, 210 checker only
and 0 validation only. Thirty of the 31 checker-only mutants remove or narrow the barrier between an \code{atomicAdd}
into a histogram and the dispatch that reads it, a read-after-write race (an atomic read-modify-write is a write) that
validation misses, confirmed by a minimal reproduction; the last narrows a readback transition to \code{dstStageMask =
ALL\_GRAPHICS}, which does not contain the copy that follows. No mutant that the checker accepted changed the output
hashes. The fresh-study-2 snapshot gives byte-identical per-mutant diagnostics, and the release gives identical
per-mutant diagnostics (95.9\%: 86 both safe / 633 both hazard / 31 checker only / 0 validation only).

Earlier, the same plans ran under validation 1.3.275 (Ubuntu package, lavapipe) and the early checker (dev). In
21 mutants (230 mutant--resource pairs) validation found hazards the checker accepted; bugs A and B explain all of them.
After the fix no mutant is validation-only: 2,199 resource pairs agree and 621 are checker-only, of which 199 are races
on storage buffers written by compute shaders that validation 1.3.275 classifies as reads, one is a copy after a barrier
whose destination is \code{ALL\_GRAPHICS}, and three are conservative checker reports. With R-LT the verdicts per mutant
and per resource are unchanged (2,199 both, 627 checker-only). Readback hashes were identical in all 750 mutants,
including those with confirmed hazards, because lavapipe executes nearly in order.

\paragraph{Ray tracing and mesh shading}
Thirteen raw-Vulkan programs on one queue: 12 equal verdicts; in the thirteenth a top-level build reads a bottom-level
structure whose build is not made available to it, and only \sys reports it. Two programs test R-LT: two transitions with
no dependency between them are clean under both tools; a clear followed by two transitions and a copy, ordered only
through the transitions, is reported by \sys, since the specification forbids relying on transition order. All 15 modes
of this test give the same codes under the engine-study-3 snapshot and every configuration of the first snapshot of
the verified configuration (post-hoc), and the release gives the same codes and counts (paper, \S\ref{P-sec:eval-sem}).
On the sample collections,
Table~\ref{tab:s-rt} gives the triage of the ray-tracing, mesh-shading and compute samples.

\begin{table}[htbp]
\caption{Ray-tracing (RT), mesh-shading and compute samples: raw reports per 10 frames (``base'': checker before ray
tracing on the same trace), triage, and validation (HEAD-max) messages. n/a: lavapipe lacks the extension.}
\label{tab:s-rt}
\scriptsize
\resizebox{\columnwidth}{!}{%
\begin{tabular}{@{}lrrr@{}}
\toprule
Samples & raw (base) & true / tech. / FA & validation \\
\midrule
SW RT, 2 samples & 48 (0) & 48 / 0 / 0 & 19 \\
SW RT, 7 samples & 120 (72) & 48 / 72 / 0 & 19 \\
SW \code{rayquery} & 288 (288) & 0 / 288 / 0 & 0 \\
SW \code{meshshader} & 0 (0) & 0 / 0 / 0 & 0 \\
KS RT, 3 samples & 94 (72) & 33 / 61 / 0 & 8 \\
KS RT, 2 samples & 10 & 0 / 10 / 0 & 0 or n/a \\
KS ray query, 2 samples & 10 (10) & 0 / 10 / 0 & 0 \\
KS mesh shading, 2 samples & 10 (10) & 0 / 10 / 0 & 0 \\
\bottomrule
\end{tabular}}
\end{table}

\subsection{Whole-frame differential test}
\label{sec:eval-wholeframe}
The checker frozen before the test agreed with validation on 88.4\% of the 1,251 programs on hazards and 84.3\% on
all finding classes (fresh). The disagreements exposed six checker defects (\S E). After they were fixed, we
regenerated all 1,251 programs on the model driver and re-checked them with the release's predecessor (\prevhash; its results are the release's, Table~\ref{tab:s-versions}),
reusing validation's stored verdicts (Table~\ref{tab:s-wholeframe}): agreement is 97.9\% on hazards and 95.4\% on all
classes; hazard pairs are 1,440 / 32 / 14 and all finding pairs 2,353 / 478 / 116 (both / \sys only / validation only).
With the \code{acquire-first-use} warning scored as a hazard, agreement is 98.3\% / 95.8\%. All 1,251 programs are
guarantee-eligible (checking complete, no model gap, no pair relying on the host-domain premise). Against the stored
results of the checker with the test's six fixes, no verdict changes in any program. The workloads have no
framebuffer-local flag, no render pass, no descriptor array and no device write followed by a host write, so the
changes made after both fresh studies have nothing to act on here; the snapshots from the first verified configuration
to the release give byte-identical diagnostics on these traces. The traces are released with the code.

The 205 host-write races are the same under both capture modes. Validation on the multi-queue baseline reports four
write-after-write false positives through a timeline round trip that \sys accepts, so that configuration is compared
relative to its baseline. On the mock driver, validation gives the same findings as on lavapipe in 967 of the 989
completed single-queue runs, all 22 differences coming from the acquire order. Speed (original run): on these programs
(six frames at 64$\times$64) a lavapipe run with validation takes a median 0.081\,s and \sys 0.092\,s, of which the
model-driver run is 0.013\,s and loading and checking 0.015\,s; at 600 frames \sys's load and check (about
0.5--0.6\,s) is slower than validation's overhead (0.1--0.23\,s). We claim no speed advantage.

\begin{table*}[t]
\caption{Whole-frame differential test, 1,251 programs. Pairs: both / \sys only / validation only. The release rows
re-check regenerated programs. The engine-study-2 and engine-study-3 snapshots (dev) and every configuration of the first
verified-configuration snapshot (post-hoc) give the same agreement and pairs as the release.}
\label{tab:s-wholeframe}
\small
\begin{tabular}{@{}lllrrll@{}}
\toprule
Checker & provenance & first-acquire counted & hazard agr. & all-class agr. & hazard pairs & all pairs \\
\midrule
snapshot frozen before the test & fresh & -- & 88.4\% & 84.3\% & 1,132 / 17 / 322 & 1,270 / 413 / 1,199 \\
release (\relhash) & post-hoc & no & 97.9\% & 95.4\% & 1,440 / 32 / 14 & 2,353 / 478 / 116 \\
release (\relhash) & post-hoc & yes & 98.3\% & 95.8\% & 1,454 / 65 / 0 & 2,367 / 511 / 102 \\
\bottomrule
\end{tabular}
\end{table*}

\paragraph{Window measurements}
Of the 187 stored traces of the two sample collections (99 Sascha Willems, 88 Khronos), 138 are periodic with a fence
bound ($N = 1$ for 26, 2 for 85, 3 for 27); 44 have no present, 4 are not periodic and 1 has no fence bound. On the 138,
checking only the first $w + N + 2$ frames found 1,229 of the 1,262 normalized diagnostic sites of the full traces while
covering 78\% of the operations; all 33 missing sites are \code{unmodelled-command} warnings from stubbed commands, and
no race or protocol site was missing. These are measurements, not instances of Theorem~\ref{P-thm:window}.

\subsection{Engine and tutorial precision, all rows}
Table~\ref{tab:s-precision} gives every precision row, fresh, dev and post-hoc, and Table~\ref{tab:s-units} the same
triage by distinct cause and by raw diagnostic, which do not depend on how report groups are merged.
\begin{table*}[t]
\caption{Timing report groups true / technicality / false alarm (/ unresolved), every snapshot. ``Frozen'': the snapshot
frozen for the study; ``E2'', ``E3'': the engine-study-2 and engine-study-3 snapshots; ``verified'' and ``observed'':
the first snapshot of the verified configuration, in the verified configuration and with observed capture. Post-hoc
rows re-use the earlier triage for every group whose hazard set is unchanged.}
\label{tab:s-precision}
\small
\begin{tabular}{@{}llllrr@{}}
\toprule
Study & snapshot & provenance & groups & true & actionable \\
\midrule
Tutorials (43) & frozen & \textbf{fresh} & 11 / 13 / 1 of 25 & 44.0\% & 96.0\% \\
 & E2 = E3 & dev & 6 / 13 / 0 of 19 & 31.6\% & 100\% \\
 & verified & post-hoc & 6 / 13 / 24 / 0 of 43 & 14.0\% & 44.2\% \\
 & observed & post-hoc & 6 / 13 / 0 / 0 of 19 & 31.6\% & 100\% \\
vkQuake + Diligent (4) & frozen & \textbf{fresh} & 8 / 0 / 44 of 52 & 15.4\% & 15.4\% \\
 & E2 = E3 & dev & 7 / 2 / 11 of 20 & 35.0\% & 45.0\% \\
 & verified & post-hoc & 7 / 2 / 34 / 0 of 43 & 16.3\% & 20.9\% \\
 & observed & post-hoc & 7 / 2 / 11 / 0 of 20 & 35.0\% & 45.0\% \\
SDL3 + bgfx (12) & frozen & \textbf{fresh} & 2 / 0 / 13 of 15 & 13.3\% & 13.3\% \\
 & E3 & dev & 2 / 0 / 0 of 2 & 100\% & 100\% \\
 & every configuration & post-hoc & 2 / 0 / 0 of 2 & 100\% & 100\% \\
Granite + VSG (12) & frozen (9 checked) & \textbf{fresh} & 8 / 12 / 0 of 20 & 40.0\% & 100\% \\
 & verified (12 checked) & post-hoc & 8 / 13 / 0 / 1 of 22 & 36.4\% & 95.5\% \\
 & heuristic R1 & post-hoc & 8 / 12 / 0 / 1 of 21 & 38.1\% & -- \\
\bottomrule
\end{tabular}
\end{table*}

\begin{table*}[t]
\caption{Merge-independent units: distinct causes and raw diagnostics, true / technicality / false alarm. Snapshot
names as in Table~\ref{tab:s-precision}.}
\label{tab:s-units}
\small
\begin{tabular}{@{}llrrrr@{}}
\toprule
Study & checker & causes & per-cause true & raw & per-raw true \\
\midrule
Tutorials & frozen (fresh) & 3 / 2 / 1 & 50.0\% & 3,208 / 63 / 70 & 96.0\% \\
Tutorials & E2 = E3 (dev) & 3 / 2 / 0 & 60.0\% & 3,208 / 63 / 0 & 98.1\% \\
vkQuake + Diligent & frozen (fresh) & 3 / 0 / 7 (+1 DRV) & 30.0\% & 84,424 / 0 / 1,747 & 97.9\% \\
vkQuake + Diligent & E2 = E3 (dev) & 3 / 1 / 2 & 50.0\% & -- & -- \\
SDL3 + bgfx & frozen (fresh) & 1 / 0 / 1 & 50.0\% & 166 / 0 / 124 & 57.2\% \\
SDL3 + bgfx & E3 (dev) & 1 / 0 / 0 & 100\% & 166 / 0 / 0 & 100\% \\
Granite + VSG & frozen (fresh) & 5 / 1 / 0 & 83.3\% & 334 / 516 / 0 & 39.3\% \\
\bottomrule
\end{tabular}
\end{table*}

\subsection{The verified configuration on every earlier corpus (post-hoc)}
Table~\ref{tab:s-regress} compares the engine-study-3 snapshot (the snapshot of the dossier counts) with the first
snapshot of the verified configuration on every earlier corpus, and Table~\ref{tab:s-engcfg} gives the engines under
every configuration on the same newly recorded traces. The release's results on the archived traces are in the paper
(Table~\ref{P-tab:release}).
\begin{table*}[t]
\caption{Earlier corpora: the engine-study-3 snapshot against the first snapshot of the verified configuration
(post-hoc), same traces where both were run.}
\label{tab:s-regress}
\small
\begin{tabular}{@{}p{0.25\textwidth}p{0.14\textwidth}p{0.14\textwidth}p{0.4\textwidth}@{}}
\toprule
Corpus & engine-study-3 snapshot & verified configuration\posthoc & why \\
\midrule
\bench strict (114) & 80 / 2 / 2 / 30 & 78 / 4 / 2 / 30 & 65, 66, 68 fixed-not-clean (true harness race, conservative capture); 94 detected (entry-point stages) \\
\bench blind timing, strict / root cause & 29 / 33 of 35 & 29 / 33 of 35 & same traces; as run: 21 / 24 \\
whole-frame (1,251) & 97.9\% / 95.4\% & 97.9\% / 95.4\% & 0 programs change \\
RT / mesh (15) & 15 / 15 & 15 / 15 & texts only \\
stored root causes & 24 IDs + T17 & 24 IDs + T17 & 0 lost (by site) \\
tutorials raw / groups & 3,291 / 22 & 9,606 / 58 & conservative capture, page-neighbour false alarms \\
SDL3 + bgfx raw / groups & 221 / 21 & 221 / 21 & -- \\
Granite raw / groups / tech. & 515 / 6 / 6 & 589 / 4 / 8 & strict \code{UNDEFINED} technicalities and P3 warning \\
VSG raw / groups / tech. & 334 / 8 / 0 (3 not checked) & 471 / 9 / 1 & CHK-PS fixed; \code{vsg\_depthpeel} final-layout technicality (\S D); message specificity \\
Diligent raw / groups & 118 / 4 & 136 / 14 & texts (+8 raw), conservative capture (+10 false-alarm groups) \\
vkQuake raw / groups & 71,101 / 14 & 81,161 / 27 & conservative capture (+10,057 raw, 13 false-alarm groups) and 3 P3 warnings \\
Sascha T1--T6 host-race raw & 502 & 1,607 & T2--T5 detected (same-value stores) \\
\bottomrule
\end{tabular}
\end{table*}

\emph{Strict transitions from \code{UNDEFINED}.} Strict mode changes no \bench status or code set, no whole-frame
verdict and no ray-tracing code, and keeps every labelled root-cause site of the stored samples. It adds diagnostics
of one shape, a write that is ordered before a transition from \code{UNDEFINED} but not made available to it, mostly
after a \code{DONT\_CARE} store; \code{report\_groups} grades them as technicalities (\S D). It merges two error groups
into one in four samples, because one \code{srcAccessMask} repair now covers both.

\emph{\bench.} Conservative capture turns three detected meaning cases (65, 66, 68) into ``fixed variant not clean'',
because both variants share a true harness race: two frames in flight share one uniform buffer, and the host rewrites
the unchanged camera after waiting only its own frame's fence. The snapshot frozen for round 4 reported this race; the
observed capture of the tutorial-study and engine-study snapshots lost it. Entry-point stage attribution turns case 94 into a clean detection.

\emph{Engines.} With the push-constant crash fixed, all 12 Granite and VulkanSceneGraph programs run:
\code{vsg\_gltf} and \code{vsg\_shadow} report nothing, and \code{vsg\_depthpeel} reports one technicality group and
one group we leave unresolved, in the report direction. Its window render pass ends with the implicit dependency, whose
second scope is empty, and a second render pass in the same command buffer then loads and writes the swapchain image.
On VSG-5, validation reports 19 cross-frame write-after-write pairs that neither snapshot reports; \sys finds each pair
ordered by a chain through the window render pass's external depth dependency into the next frame's dependency, which
the specification text supports, and removing those dependencies from the model makes strict mode report them
(\S\ref{sup:vsg5}).

\begin{table*}[t]
\caption{Engines on newly recorded traces, every configuration of the first snapshot of the verified configuration
(post-hoc), and the engine-study-3 snapshot on the same traces: raw timing / report groups / technicality groups.}
\label{tab:s-engcfg}
\small
\begin{tabular}{@{}llllll@{}}
\toprule
Programs & engine-study-3 snapshot & \textbf{verified} & observed capture & heuristic R1 & R1, observed \\
\midrule
SDL3 + bgfx (12) & 221/21/0 & 221/21/0 & 221/21/0 & 221/21/0 & 221/21/0 \\
Granite (4) & 515/6/6 & 589/4/8 & 589/4/8 & 519/6/6 & 519/6/6 \\
VSG (8) & 334/8/0$^*$ & 471/9/1 & 471/9/1 & 431/9/0 & 431/9/0 \\
Diligent & 118/4/2 & 136/14/2 & 126/4/2 & 136/14/2 & 126/4/2 \\
vkQuake (3) & 71,101/14/0 & 81,161/27/0 & 71,104/14/0 & 81,161/27/0 & 71,104/14/0 \\
\bottomrule
\multicolumn{6}{@{}l}{$^*$3 programs not checked: the check aborted (CHK-PS).}
\end{tabular}
\end{table*}

\subsection{\rcgp}
\label{sup:rcgp}
For every first-collection bug we asked whether \rcgp's contract language rules it out through its typed API (the buggy
program fails to compile or cannot be written; escape hatches do not count), from its API and source (Table~\ref{tab:s-software}), confirmed by 35 compile-only ports compiled by its own compiler, each buggy and fixed; where a port
disagreed with our reading, the compile result was kept. \rcgp can express 23 of 126: 22 of 27 structure bugs, 1 of 48
timing bugs and none of 51 meaning bugs. The ports changed seven earlier verdicts. Head to head, we ported every blind
case \rcgp can express plus the 12 closest calls (29 ports, 58 compiles, all as predicted): \rcgp rules out 17 of the 98
blind cases (structure 16 of 21, timing 1 of 35, meaning 0 of 42), \sys detects 29 under strict scoring; both catch 6,
\rcgp alone 11 (all structure), \sys alone 23, either 40. \rcgp's four public samples under the tutorial-study snapshot give no report and no false
alarm; with a legal \code{VK\_SUBOPTIMAL\_KHR} from acquire, \sys reports a swapchain over-acquire in every sample, caused
by \rcgp's runtime treating a suboptimal acquire as failed (this rule was written after we saw a sample freeze, so it is
not independent evidence).

\subsection{Barrier transformations}
\label{sec:eval-elide}
These measurements were made with an early optimizer and checker (development data) and were not re-run with the
release (\S\ref{sec:elide}).

\paragraph{Microbenchmarks}
Three headless workloads we wrote with a full barrier after every pass: \emph{particles} (8 independent compute
dispatches, binning with atomics, shading), \emph{deferred} (three G-buffer passes, lighting, tonemap) and
\emph{postchain} (two interleaved post-processing chains of six levels). The plans remove 63 of 111 recorded barrier
calls, narrow 45 and keep 3. On lavapipe the output hashes are identical and validation reports zero hazards under
strict matching (3 traced frames) and chain matching (8 frames). A naive control that removes every barrier redundant
on its own creates 30 hazards on \emph{postchain}, and the checker and validation each report 30.
Table~\ref{tab:s-speedup} reports GPU time on a consumer GPU: 15 runs of 300 frames per configuration, first 30 dropped,
timestamp queries, bootstrap 95\% intervals. The eight-frame chain-matched runs are measurements, not instances of a
guarantee (\S\ref{sec:elide}).

\begin{table*}[t]
\caption{Consumer GPU: median GPU frame time (ms) and GPU time saved (95\% CI).}
\label{tab:s-speedup}
\small
\begin{tabular}{@{}lrrrrl@{}}
\toprule
Workload & scale & base & layer, off & plan & saved \\
\midrule
particles & 1 & 0.236 & 0.233 & 0.197 & 16.7\% (15.6..17.0) \\
 & 4 & 0.800 & 0.800 & 0.752 & 5.9\% (5.8..6.1) \\
 & 16 & 3.310 & 3.311 & 3.258 & 1.6\% ($-$1.3..4.7) \\
deferred & 1 & 0.077 & 0.077 & 0.058 & 24.4\% (23.9..24.6) \\
 & 4 & 0.230 & 0.230 & 0.211 & 8.4\% (8.2..8.6) \\
 & 16 & 0.886 & 0.882 & 0.867 & 2.2\% (0.8..7.6) \\
postchain & 1 & 0.140 & 0.141 & 0.107 & 23.6\% (22.6..24.5) \\
 & 4 & 0.420 & 0.427 & 0.367 & 12.5\% (11.6..14.5) \\
 & 16 & 1.708 & 1.630 & 1.655 & 3.1\% ($-$1.8..7.9) \\
\bottomrule
\end{tabular}
\end{table*}

\paragraph{Real samples: a null result}
On eight unmodified samples (three Sascha Willems, five Khronos) the prover removes or narrows none of 414 per-frame
barrier calls: 318 sit between accesses that already race and are pinned, and the other 96 are transitions it proved
necessary. Only 150 one-time setup barriers are narrowed. On lavapipe the presented-image hashes are identical between
base and plan, and validation's messages do not change.

\subsection{Cost}
The checks are offline. The times below were measured with the snapshot in use for each corpus when it was run, not
with the release; the release's cost on the largest fresh-study programs is in the paper. With the epoch mode, the timing check of the slowest Sascha Willems sample takes 5.1\,s; all
Sascha Willems samples together take 44.7\,s under three-way parallel load (about 17\,s idle), and Khronos 302\,s, its four
largest traces (over 400k records) taking 50--74\,s each, mostly loading. Grouping, including the oracle runs that
re-check every candidate repair, costs 136\,s on Sascha Willems and 285\,s on Khronos. Checking the nine Granite and
VulkanSceneGraph programs took 78\,s and grouping 79\,s (engine-study-3 snapshot); vkQuake's timedemos took
105--271\,s to check and 331--650\,s to group (snapshot frozen for engine study 1); the 43 tutorial traces 182\,s and 422\,s. Table~\ref{tab:s-overhead} breaks the cost down on eight
samples measured one at a time on a 4-core machine with no GPU (an early checker around the tutorial-study snapshot):
tracing is cheaper than running the
same 10 frames on lavapipe in every case, because the model driver executes no GPU work, and the timing check takes at
most 1.6\,s.

\begin{table*}[t]
\caption{Overhead on eight samples (10 frames, early checker around the tutorial-study snapshot, measured one at a
time). Trace and lavapipe: wall time
with a warm file cache. Check phases in seconds; grouping includes its oracle runs. Peak resident memory of the traced
application and of the same application on lavapipe.}
\label{tab:s-overhead}
\small
\begin{tabular}{@{}lrrrrrrrrr@{}}
\toprule
& & & \multicolumn{2}{c}{run (s)} & \multicolumn{4}{c}{check phases (s)} & \\
\cmidrule(lr){4-5}\cmidrule(lr){6-9}
Sample & records & trace MB & traced & lavapipe & load & timing & struct. & grouping & peak MB (app / lavapipe) \\
\midrule
sw/triangle & 318 & 0.04 & 0.005 & 0.13 & 0.04 & 0.02 & 0.00 & 0.01 & 10 / 100 \\
sw/deferred & 719 & 3.3 & 0.17 & 1.22 & 0.10 & 0.09 & 0.01 & 0.77 & 25 / 255 \\
sw/computecloth & 2,545 & 1.9 & 0.06 & 0.27 & 0.10 & 0.10 & 0.28 & 0.21 & 23 / 118 \\
sw/oit & 3,122 & 1.1 & 0.03 & 0.95 & 0.19 & 1.49 & 0.07 & 11.36 & 10 / 184 \\
sw/shadowmappingomni & 2,901 & 1.1 & 0.03 & 0.38 & 0.17 & 0.74 & 0.07 & 15.38 & 10 / 139 \\
sw/multithreading & 31,678 & 6.0 & 0.28 & 2.25 & 0.37 & 1.61 & 0.31 & 1.39 & 12 / 192 \\
ks/hdr & 644 & 2.3 & 0.21 & 1.64 & 0.50 & 0.11 & 0.01 & 0.10 & 141 / 202 \\
ks/afbc & 5,404 & 34.7 & 2.35 & 3.53 & 0.36 & 0.22 & 0.04 & 0.22 & 145 / 496 \\
\bottomrule
\end{tabular}
\end{table*}

%% ===================================================================
\section{The Two Fresh Studies of the Proved Configuration}
\label{sup:fresh}

This section gives the protocols, triage rules and full results of the two studies of the paper's E3. Every number
is from the checker snapshot frozen for its study, in the configuration named, except in \S\ref{sup:release-posthoc},
which re-checks the same recorded traces with the release and is post-hoc. The artifact holds the traces, the
per-diagnostic classes and the scripts that produce each table.

\subsection{Protocols and disclosures}
\emph{Order of events.} For each study: (1) candidate programs were built and brought up under the model driver,
recording only whether a program starts and how many calls, presents, draws and host writes it produces; no check,
grouping or validation run was made on any bring-up trace; (2) the protocol fixing the program set, the run length (20
presents), the commands, the triage rules and the metrics was written; (3) the checker was frozen by source hash and
its full test suite run from the frozen snapshot; (4) every program was traced, checked in both capture modes,
grouped, and run under validation; (5) report groups were triaged by the written rules. Amendments to a protocol are
timestamped and append-only.

\emph{Study 1}: 32 programs. Qt~6.4.2 rendering hardware interface with its Vulkan
backend, 13 Qt Quick~3D examples and one Qt Quick example, run with the \code{qml} tool; wgpu (15 examples, Vulkan
backend, release build); the two Vulkan examples of Dear ImGui; and vkcube~1.3.275. Programs that render only on
demand (fewer than 20 presents), need a compiled module, or redraw only on input were excluded at bring-up. Prior
exposure: Dear ImGui's Vulkan backend had been linked into two earlier programs as an overlay; no rule was written
for it. Amendment: the sequential runner was replaced by a three-wide one. Bring-up found one model-driver defect,
fixed before the freeze; it changed no checker logic.

\emph{Study 2}: 25 programs. GTK~4.14.5 with its Vulkan renderer
(\code{gtk4-widget-factory} and eight \code{gtk4-demo} demos: fishbowl, paintable\_animated, iconscroll,
path\_spinner, path\_walk, spinner, mask, listview\_clocks); vkmark (scenes vertex, texture, shading, effect2d,
desktop, cube, clear); mpv~0.37 with libplacebo~6.338 (\code{gpu} and \code{gpu-next}, default and high-quality
profiles) on a generated test video; and five GStreamer~1.24.2 Vulkan pipelines (upload to sink, colour conversion,
overlay composition, view conversion, a SPIR-V shader). Seven GTK demos were excluded at bring-up (fewer than 20
presents). Study~2's snapshot is study~1's plus the fixes of the three defects study~1 exposed and two commands
bring-up showed were needed (\code{vkCmdClearAttachments}, \code{vkCmdCopyQueryPoolResults}), each tested on the model driver's own
example, never on a study program. Prior exposure: a libplacebo merge request had been considered and rejected as an
earlier benchmark case; no libplacebo code had been run. Amendments: (a) the first start traced only one program
because the program table returned a wrong status; no check output was looked at, the run directory was deleted and
the study restarted; (b) the conservative check of \code{gtk\_path\_walk} was killed by the container's memory limit
while other checks ran beside it and was rerun alone; it then finished, but grouping exceeded the memory limit even
alone, so the program is checked but not grouped and its diagnostics were classified one by one; (c) libplacebo refuses software devices unless one is named, so
the mpv baseline runs name lavapipe.

\emph{Baseline.} Khronos validation layer (Table~\ref{tab:s-software}) with synchronization validation and its shader-access
heuristic and load-after-store checks enabled, on lavapipe (Mesa~25.2.8), same programs and commands, limited to 20
presents by a shim (mpv: to the end of the test video, since libplacebo loads its own entry points).

\subsection{Triage rules}
The rules were written before study~1's first check. \emph{Families}: a \emph{race} diagnostic reports two
unordered accesses (layout transitions and host writes included); a \emph{protocol} diagnostic reports a broken
valid-usage rule of a synchronization object (a binary semaphore signalled while a wait may be pending, a wait with
no possible signal, a fence or command buffer reused while in use, swapchain over-acquire, timeline value order).
\emph{Classes}: \emph{true} (a concrete rule is broken and either a schedule the specification allows changes what
the program reads or presents, or the rule is a valid-usage rule); \emph{technicality} (a race whose rule is broken
but no schedule changes what is read or presented, in exactly five named sub-classes: a transition from
\code{UNDEFINED} against discarded contents, a layout-name mismatch with the same memory layout, a missing ownership
transfer within one family, ordering only through the ordering of layout transitions among themselves, and a store of
bytes already present); \emph{false alarm} (the program is correct; the cause is named); \emph{driver} (the model
driver's fault; excluded); and \emph{unresolved} (needs a ruling on the specification; counted against precision).
Semaphore reuse that relies on a queue idle or a fence wait is not a class of its own: every protocol violation is
true. \emph{Amendments}: a report group takes the most consequential class of its members
(true, unresolved, technicality, false alarm, driver); a checker defect found during triage is not fixed in the frozen
checker. Each program also gets a \emph{premise verdict}: ``no detected violation'' when there is no capture warning (P1,
P3), no unmodelled command and no untracked mechanism; otherwise outside, with the warnings named. ``No detected
violation'' is not a proof that the premises hold: P2 and the model premises are assumed, and P3 is checked only at the
first store to each page.

\subsection{Study 1: causes}
\begin{table*}[t]
\caption{Study 1: causes. Groups in conservative capture. Upstream: the current source of the project.}
\label{tab:s-fresh1}
\scriptsize
\begin{tabular}{@{}p{0.43\textwidth}llp{0.17\textwidth}rp{0.14\textwidth}@{}}
\toprule
Cause & class & family & programs & groups & validation / upstream \\
\midrule
QT-ACQ acquire semaphore reused with a pending wait & true & protocol & 14 Qt & 56 & 01779 / fixed \\
QT-PSR present semaphore per frame slot, not per image & true & protocol & 14 Qt & 28 & 00067 / fixed \\
VKCUBE-PSR present semaphore per frame & true & protocol & vkcube & 2 & 00067 / fixed \\
IMGUI-PRE swapchain images transitioned before they are acquired & true & race & 2 ImGui & 40 & non-acquired image / \textbf{present} \\
QT-VS texture made visible to the vertex stage only, sampled in the fragment shader & true & race & 4 Qt & 4 & RAW / \textbf{present} \\
QT-L, QT-LF render-pass \code{UNDEFINED} transition right after Qt's own barrier & tech. (L) & race & 10 Qt & 12 & silent \\
QT-U multisampled colour buffer shared by consecutive frames & tech. (U) & race & 2 Qt & (in QT-L) & WAW \\
FA-PG conservative page capture of per-slot uniform writes & false alarm & race & 14 Qt (cons.\ only) & 1,026 & silent \\
FA-CQ timeline counter query not taken as a host wait & false alarm & both & 15 wgpu & 934 & silent / fixed for study 2 \\
FA-WL whole-image layout as an unbounded range & false alarm & race & 2 wgpu & 10 & silent / fixed for study 2 \\
FA-X8 \code{X8\_D24} treated as colour & false alarm & race & 1 wgpu & 1 & silent / fixed for study 2 \\
\bottomrule
\end{tabular}
\end{table*}

Dear ImGui's 20 groups per program come from the examples recreating the swapchain every frame under the model
driver's fixed surface extent; they are one cause. Re-checking the 15 wgpu programs with the three defects fixed
(post-hoc) leaves all 15 clean in both modes; study~1 as a whole is then 130 / 12 / 1,026 groups in conservative and
130 / 12 / 0 in observed capture. Validation found the same five true causes; it is silent on QT-L and QT-LF and on
every false alarm. One wgpu program crashed under lavapipe, so validation has no result for it.

\subsection{Study 2: causes}
\begin{table*}[t]
\caption{Study 2: causes. Groups in conservative capture (observed capture is identical). Diagnostics: report-group
members. Upstream: the current source of the project.}
\label{tab:s-fresh2}
\scriptsize
\begin{tabular}{@{}p{0.47\textwidth}llp{0.13\textwidth}rp{0.15\textwidth}@{}}
\toprule
Cause & class & family & programs & groups & validation / upstream \\
\midrule
VKMARK-PSR one submit semaphore per scene, waited by every present & true & protocol & 7 vkmark & 7 & 00067 / \textbf{present} \\
VKMARK-ACQ one acquire semaphore, no fence before the next acquire & true & protocol & 7 vkmark & 16 & 01779 / fixed \\
VKMARK-UBO uniform buffer rewritten by \code{memcpy} every frame, no wait for the frame reading it & true & race & 5 vkmark & 8 & \textbf{silent} / fixed \\
VKMARK-CLR clear scene: barrier source stage \code{TRANSFER} does not chain to the acquire wait at \code{COLOR\_ATTACHMENT\_OUTPUT} & true & race & clear & 4 & WAR / \textbf{present} \\
VKMARK-DEPTH one depth image for all frames, render pass dependency on colour output only & true & race & 3 vkmark & 3 & WAW / \textbf{present} \\
GST-WAIT acquire semaphore waited at \code{BOTTOM\_OF\_PIPE} & true & race & 5 gst & 30 & WAR / \textbf{present} \\
GST-IA sampled image made visible to input-attachment reads only & true & race & 3 gst & 3 & RAW / \textbf{present} \\
PLACEBO-EXT first barrier after an external semaphore wait has source stage \code{NONE} & true & race & 4 mpv & 16 & WAR / fixed \\
GTK-SEM one acquire semaphore, re-signalled while the present's wait may pend & true & protocol & 8 gtk & 8 & 01779 / fixed \\
GTK-ACQ rendering never waits for the acquire semaphore & true & race & 8 gtk & 47 & non-acquired image, WAP / fixed \\
GTK-PRS present waits for the acquire semaphore, not the rendering & true & race & 8 gtk & 118 & PAW / fixed \\
VKMARK-DEPTH-U the render pass's \code{UNDEFINED} transition against the previous frame's depth writes & tech. (U) & race & 3 vkmark & (mixed) & in WAW \\
GTK-TL barrier transition, then the render pass's own transition & tech. (L) & race & iconscroll & (mixed) & silent \\
GTK-DESC read charged to descriptor-array elements the draw does not sample & false alarm & race & 4 gtk & 1,042 & silent \\
GTK-ATLAS uploads into disjoint rectangles of one atlas & false alarm & race & 6 gtk & 6 & silent \\
GTK-RPX render-pass final transition vs.\ a later \code{COLOR\_ATTACHMENT\_OUTPUT} barrier & unresolved & race & 4 gtk & 7 & silent \\
GTK-MP render passes per damage rectangle, no barrier between & unresolved & race & 4 gtk & 6 & silent \\
\bottomrule
\end{tabular}
\end{table*}

\emph{GTK-DESC, evidence.} GTK's renderer encodes each image's descriptor-array index in a per-instance vertex
attribute and samples \code{textures[nonuniformEXT(id >> 1)]}. We rebuilt every host-visible allocation at each
submit from the recorded host writes, read the instance records of every draw through its pipeline's vertex input
state and the vertex shader's attribute locations, and mapped each index to an image through the last descriptor
update. All 11,875 decoded indices fall on written elements. 135,709 diagnostics name an image their draw does not
sample, including every ``wrong layout for a shader read'' diagnostic; only 395 name a sampled image, all of them
GTK-RPX. The specification counts only dynamically used descriptors as accessed, so these are false alarms. A static
model of the recorded trace could decode the indices the same way; \sys does not.

\emph{GTK-RPX and GTK-MP.} GTK's render passes declare no dependencies and so get only the implicit external ones,
whose exit dependency has an empty second scope. GTK-RPX: a later barrier with source stage
\code{COLOR\_ATTACHMENT\_OUTPUT} makes the pass's colour writes visible to a sampled read; whether its first scope also
contains the pass's final layout transition is not stated. GTK-MP: several render passes on one swapchain image per
frame, one per damage rectangle, with no barrier between them; whether two render passes' automatic transitions are
ordered with each other is not stated (the specification orders transitions of image barriers with other transitions,
including those of render passes, and says that ordering must not be relied on). Validation accepts both. They are
one question, shared with \code{vsg\_depthpeel}, and need a ruling. Both lie in the report direction: \sys reports the
pairs and validation accepts them, so no clean verdict depends on them (\S\ref{sup:vsg5}).

\emph{Validation.} The tools share 10 true causes; VKMARK-UBO, a host-write race, is unique to \sys; none is unique to
validation. Validation reports out-of-scope messages in vkmark (\code{VUID-VkShaderModuleCreateInfo-pCode-08740},
\code{VUID-VkImageCreateInfo-initialLayout-12478}).

\subsection{Provenance of upstream fixes}
\label{sup:ledger}
Both studies ran after every fix below landed, and we reported nothing upstream before or since; all nine commits
predate the studies, and none was caused by this work. Fixing commits were found with \code{git log -S/-G} on the lines
that carry each cause, and a commit is listed only when its diff changes the exact pattern. Days: from the commit's
landing on the default branch to our study. The GTK fixes landed on main before the tested 4.14.5 release but were not
backported to the 4.14 branch. The full ledger, with release tags and backports, is in the artifact.

Each commit carries one of three labels. An \emph{initial fix attempt} was re-run under the same commands and leaves
diagnostics of the cause. A \emph{commit changing the pattern} was read in the source and not re-run. A \emph{verified
complete repair} would be re-run with no diagnostic of the cause left; none of the nine was re-run as a complete
repair.

\begin{table}[h]
\caption{The nine true causes with an upstream commit that changes the implicated pattern. Label: initial fix attempt
(re-run, leaves diagnostics of the cause) or commit changing the pattern (read, not re-run).}
\label{tab:s-provenance}
\scriptsize
\resizebox{\columnwidth}{!}{%
\begin{tabular}{@{}lllllr@{}}
\toprule
cause & tested rev. & upstream commit & label & issue named & days \\
\midrule
QT-ACQ & Qt 6.4.2 & \code{24d62ffd720} (2025-01-27) & changes the pattern & QTBUG-132808 & 627 \\
QT-PSR & Qt 6.4.2 & \code{eab1295c882} (2025-06-25) & changes the pattern & QTBUG-137986 & 478 \\
VKCUBE-PSR & 1.3.275 & \code{a8f20787} (2025-04-11) & changes the pattern & -- & 553 \\
VKMARK-ACQ & \code{ab6e6f3} & \code{956e145} (2025-01-08) & initial fix attempt & -- & 646 \\
VKMARK-UBO & \code{ab6e6f3} & \code{fd2d798} (2025-02-25) & initial fix attempt & vkmark\#69 & 598 \\
PLACEBO-EXT & v6.338.2 & \code{fd01c75c} (2025-03-14) & changes the pattern & libplacebo\#333 & 581 \\
GTK-SEM & 4.14.5 & \code{b08a17064e} (2024-07-14) & changes the pattern & -- & 824 \\
GTK-ACQ & 4.14.5 & \code{4966f8cdf8} (2024-07-17) & changes the pattern & -- & 821 \\
GTK-PRS & 4.14.5 & \code{fc6581cf3b} (2025-11-12) & changes the pattern & gtk\#7866 & 338 \\
\bottomrule
\end{tabular}}
\end{table}

The two initial fix attempts are vkmark's. The uniform-buffer commit (one uniform buffer per swapchain
image) leaves 136 of the cause's 152 diagnostics, because it rewrites image $i$'s buffer once acquire returns $i$, which
does not show the host that the frame that read that buffer has finished; and in the \code{shading} scene it drops the
fragment stage from the uniform binding's descriptor layout (\code{structure.binding-not-in-layout}; validation reports
VUID 07988), which vkmark corrected three weeks later. The acquire commit leaves 119
\code{acquire-semaphore-in-use} diagnostics, and validation still reports its semaphore error (01779).

Seven causes are present in the current sources (checked 2026-10-16). Two were reported by others before our study:
IMGUI-PRE (Dear ImGui issue \#9496 and open PR \#9553) and VKMARK-PSR (vkmark issue \#88 and open PR \#89, which does
not touch VKMARK-CLR or VKMARK-DEPTH). We found no report of VKMARK-CLR or VKMARK-DEPTH. For QT-VS, GST-WAIT and
GST-IA we could not search the projects' trackers.

\subsection{Repair-to-clean cases}
\label{sup:repairs}
There are \emph{11 source-patch cases covering nine causes}: one case per cause (SDL-1, IMGUI-PRE, GST-WAIT, GST-IA and
the five vkmark causes) and two combined cases, GST-ALL and VKMARK-ALL, that apply all patches of a project together.
They are distinct from the 11 historical dossiers of \S\ref{sup:dossiers}; SDL-1 is in both. Each case: the original
program, a source repair, a rebuild with the original flags, and a rerun with the same command (model driver, 20
presents), checked in conservative and observed capture, with validation (lavapipe) before and after.
A repair is \emph{clean for its cause} when the patched run has no diagnostic of that cause and no diagnostic code the
original lacked; remaining diagnostics are attributed by written rules (code and resource) to the other listed causes.
The cases were first checked with the snapshot frozen for study~2 and then re-traced from the same sources and patches
and re-checked with the release (post-hoc; Table~\ref{tab:s-repairs}). The tested vkmark package is source-identical to the commit in Table~\ref{tab:s-software}. Patches, commands, traces and data are in the artifact.

\begin{table*}[t]
\caption{The 11 source-patch cases. Diagnostics / groups, conservative capture, summed over the case's programs, with
the fresh-study-2 snapshot; observed capture gives the same in every run. Release\posthoc: diagnostics of the case's
cause (for GST-ALL and VKMARK-ALL, of all the project's causes) before $\to$ after, on re-traced runs; every case is
clean for its cause under the release, and every conservative run is guarantee-eligible.}
\label{tab:s-repairs}
\scriptsize
\begin{tabular}{@{}lllllll@{}}
\toprule
case & programs & patch & before & after (remaining causes) & release\posthoc, cause & clean \\
\midrule
SDL-1 & 1 & +38/$-$18 & 166 / 2 & 0 / 0 & 166 $\to$ 0 & yes \\
IMGUI-PRE & 2 & +0/$-$75 & 40 / 40 & 0 / 0 & 40 $\to$ 0 & yes \\
GST-WAIT & 5 & +1/$-$1 & 495 / 33 & 95 / 3 (GST-IA) & 400 $\to$ 0 & yes \\
GST-IA & 5 & +1/$-$1 & 495 / 33 & 400 / 30 (GST-WAIT) & 95 $\to$ 0 & yes \\
GST-ALL & 5 & +2/$-$2 & 495 / 33 & 0 / 0 & 495 $\to$ 0 & yes \\
VKMARK-UBO & 5 & +45/$-$5 & 978 / 26 & 370 / 18 (ACQ, DEPTH, PSR) & 152 $\to$ 0 & yes \\
VKMARK-PSR & 7 & +42/$-$35 & 1124 / 38 & 991 / 31 (ACQ, CLR, DEPTH, UBO) & 133 $\to$ 0 & yes \\
VKMARK-ACQ & 7 & +47/$-$18 & 1124 / 38 & 1028 / 30 (CLR, DEPTH, PSR, UBO) & 143 $\to$ 0 & yes \\
VKMARK-CLR & 1 & +1/$-$1 & 108 / 9 & 48 / 5 (ACQ, PSR) & 60 $\to$ 0 & yes \\
VKMARK-DEPTH & 3 & +10/$-$4 & 807 / 15 & 171 / 12 (ACQ, PSR, UBO) & 636 $\to$ 0 & yes \\
VKMARK-ALL & 7 & +135/$-$58 & 1124 / 38 & 0 / 0 & 1,124 $\to$ 0 & yes \\
\bottomrule
\end{tabular}
\end{table*}

No run had a capture warning, an unmodelled command or an unimplemented command, and no diagnostic was left
unattributed to a listed cause. With all patches of a project applied, 15 programs have no diagnostic in either capture
mode under the release and no synchronization message under validation: all seven vkmark scenes, all five GStreamer
pipelines, both ImGui examples and SDL's sprite test. Each of the 15 was traced three more times and once with the
release's own build of the model driver, with the same result. All 152 re-traced runs of the cases (originals,
single-cause variants, upstream commits and repairs) were checked with the release in both capture modes; every check
completed, and no pair relies on the host-domain premise. One run-to-run difference is not a synchronization finding: in
some runs, original or patched, \code{gst\_colorconvert} draws one \code{structure.host-layout} warning, because
GStreamer copies a uniform struct whose padding it never initializes; it is outside the repaired causes and outside
the timing check, and adds one unattributed diagnostic to the GStreamer totals of the re-traced runs (496 before
instead of 495). PLACEBO-EXT was not attempted: the mpv programs are outside the premises. Qt and GTK were not
rebuilt.

\subsection{The release on the recorded traces (post-hoc)}
\label{sup:release-posthoc}
Everything in this subsection is post-hoc and is kept apart from the frozen results above. After both studies the
checker was changed in four ways: framebuffer-local dependencies are read (\S\ref{sup:fblocal}); two automatic
render-pass transitions on one queue are no longer exempt from each other (\S\ref{sup:ltpairs}); a descriptor-array
read is narrowed to the elements its index bytes select when the guard holds (\S\ref{sec:impl}); and grouping has a
memory and pair budget. When the budget is exceeded, grouping reports \emph{incomplete}, keeps every diagnostic in a
site group and proposes no repair; a stub marked incomplete is written first, so a run that is killed never reads as
clean. The release re-checked the traces both studies recorded, with the frozen triage rules and one added rule: a
pair of automatic render-pass transitions with no dependency between them is \emph{unresolved}, the question of
GTK-MP. On these traces the release loses no true group of the frozen runs.

\emph{Study 1} (32 programs, all grouped completely): in conservative capture 130 / 10 / 2 / 1,026 groups (true /
technicality / unresolved / false alarm), 11.1\% true; in observed capture 130 / 10 / 2 / 0, 91.5\% true. Every false
alarm is Qt's page-granular host capture (FA-PG). The two unresolved groups are pairs of automatic render-pass
transitions in two Qt programs. All 15 wgpu programs report nothing, \code{wgpu\_mipmap} included.

\emph{Study 2}: the 24 programs whose grouping completes give 260 / 1 / 8 / 6 groups, so 94.5\% of these 275 groups are
true (19.7\% when frozen); observed capture gives the same. The false alarms fall from 1,048 groups to 6 because a draw
that selects its texture with a flat per-instance index whose bytes the host wrote, with no racing write, is charged
only the elements those bytes select; the six that remain are GTK-ATLAS. \code{gtk\_path\_walk} reaches the grouping
budget and is reported incomplete: 35 site groups with no repairs, of which 22 are true (the same three GTK causes), 12
unresolved and 1 a false alarm, with all 15,595 diagnostics kept. It is not in the 94.5\%.

\emph{Domain.} In conservative capture all 57 runs (32 + 25) are guarantee-eligible: checking completed, no
unmodelled command has a memory effect, and no host write raises the host-domain warning. The 57 observed-capture runs
are outside the verified configuration by definition. Reporting is complete for 56 programs; \code{gtk\_path\_walk}
is the exception.

\subsection{Per-program results}

\begin{table*}[p]
\caption{Per-program results of both studies, each with the snapshot frozen for it (fresh). Groups: true / technicality / false alarm (/ unresolved). Premise: ``in'' = no detected
violation, ``out'' = outside (study 1: \code{wgpu\_mipmap}, query result copies; study 2: mpv, timestamp queries). Check: seconds of the
timing check in conservative and observed capture.}
\label{tab:s-perprog}
\scriptsize
\begin{tabular}{@{}llllllr@{}}
\toprule
Study & program & premise & groups, conservative & groups, observed & \multicolumn{2}{r}{check (s), cons.\ / obs.} \\
\midrule
1 & \code{imgui\_glfw} & in & 20/0/0 & 20/0/0 & 0.2 & 0.2 \\
1 & \code{imgui\_sdl2} & in & 20/0/0 & 20/0/0 & 0.3 & 0.2 \\
1 & \code{q3d\_custommaterial} & in & 6/1/113 & 6/1/0 & 84.6 & 72.3 \\
1 & \code{q3d\_dynamiccreation} & in & 6/0/73 & 6/0/0 & 173.5 & 181.9 \\
1 & \code{q3d\_hellocube} & in & 7/1/45 & 7/1/0 & 24.0 & 27.1 \\
1 & \code{q3d\_instancing} & in & 6/1/16 & 6/1/0 & 21.0 & 21.2 \\
1 & \code{q3d\_intro} & in & 6/1/11 & 6/1/0 & 7.8 & 7.6 \\
1 & \code{q3d\_lights} & in & 6/2/51 & 6/2/0 & 47.7 & 47.2 \\
1 & \code{q3d\_morphing} & in & 6/1/14 & 6/1/0 & 15.5 & 15.2 \\
1 & \code{q3d\_offlineshaders} & in & 6/0/3 & 6/0/0 & 9.9 & 9.5 \\
1 & \code{q3d\_picking} & in & 7/0/22 & 7/0/0 & 69.8 & 67.0 \\
1 & \code{q3d\_quickball} & in & 7/1/93 & 7/1/0 & 18.2 & 16.4 \\
1 & \code{q3d\_reflectionprobes} & in & 6/1/396 & 6/1/0 & 61.0 & 61.8 \\
1 & \code{q3d\_screenspacereflections} & in & 7/2/81 & 7/2/0 & 27.9 & 28.5 \\
1 & \code{q3d\_view3d} & in & 6/1/72 & 6/1/0 & 85.5 & 84.6 \\
1 & \code{qq\_shadereffects} & in & 6/0/36 & 6/0/0 & 4.0 & 3.8 \\
1 & \code{vkcube} & in & 2/0/0 & 2/0/0 & 0.2 & 0.3 \\
1 & \code{wgpu\_boids} & in & 0/0/89 & 0/0/89 & 0.2 & 0.3 \\
1 & \code{wgpu\_bunnymark} & in & 0/0/63 & 0/0/63 & 0.9 & 1.0 \\
1 & \code{wgpu\_conservative\_raster} & in & 0/0/62 & 0/0/62 & 0.2 & 0.3 \\
1 & \code{wgpu\_cube} & in & 0/0/43 & 0/0/43 & 0.2 & 0.2 \\
1 & \code{wgpu\_mipmap} & out & 0/0/69 & 0/0/69 & 0.3 & 0.3 \\
1 & \code{wgpu\_msaa\_line} & in & 0/0/43 & 0/0/43 & 0.1 & 0.1 \\
1 & \code{wgpu\_multiple\_render\_targets} & in & 0/0/62 & 0/0/62 & 0.2 & 0.2 \\
1 & \code{wgpu\_render\_to\_texture} & in & 0/0/6 & 0/0/6 & 0.0 & 0.0 \\
1 & \code{wgpu\_render\_with\_compute} & in & 0/0/35 & 0/0/35 & 0.4 & 0.4 \\
1 & \code{wgpu\_shadow} & in & 0/0/227 & 0/0/193 & 0.8 & 0.7 \\
1 & \code{wgpu\_skybox} & in & 0/0/67 & 0/0/63 & 0.5 & 0.6 \\
1 & \code{wgpu\_srgb\_blend} & in & 0/0/43 & 0/0/43 & 0.1 & 0.1 \\
1 & \code{wgpu\_stencil\_triangles} & in & 0/0/43 & 0/0/43 & 0.1 & 0.1 \\
1 & \code{wgpu\_storage\_texture} & in & 0/0/6 & 0/0/6 & 0.7 & 0.8 \\
1 & \code{wgpu\_water} & in & 0/0/87 & 0/0/86 & 8.7 & 8.3 \\
2 & \code{gst\_colorconvert} & in & 7/0/0 & 7/0/0 & 1.2 & 1.3 \\
2 & \code{gst\_overlay} & in & 6/0/0 & 6/0/0 & 0.1 & 0.1 \\
2 & \code{gst\_shaderspv} & in & 7/0/0 & 7/0/0 & 0.1 & 0.1 \\
2 & \code{gst\_sink} & in & 6/0/0 & 6/0/0 & 0.1 & 0.0 \\
2 & \code{gst\_viewconvert} & in & 7/0/0 & 7/0/0 & 0.3 & 0.3 \\
2 & \code{gtk\_fishbowl} & in & 22/0/17/3 & 22/0/17/3 & 46.3 & 44.9 \\
2 & \code{gtk\_iconscroll} & in & 22/0/928/4 & 22/0/928/4 & 308.6 & 296.4 \\
2 & \code{gtk\_listview\_clocks} & in & 22/0/1 & 22/0/1 & 14.5 & 14.9 \\
2 & \code{gtk\_mask} & in & 22/0/1 & 22/0/1 & 32.0 & 34.8 \\
2 & \code{gtk\_paintable\_animated} & in & 22/0/0 & 22/0/0 & 2.4 & 2.4 \\
2 & \code{gtk\_path\_spinner} & in & 19/0/0 & 19/0/0 & 3.6 & 3.6 \\
2 & \code{gtk\_path\_walk} & in & not grouped & not grouped & 718.8 & 896.1 \\
2 & \code{gtk\_spinner} & in & 22/0/5/2 & 22/0/5/2 & 14.8 & 14.9 \\
2 & \code{gtk\_widget\_factory} & in & 22/0/96/4 & 22/0/96/4 & 104.8 & 110.4 \\
2 & \code{mpv\_gpu} & out & 4/0/0 & 4/0/0 & 0.5 & 0.5 \\
2 & \code{mpv\_gpu\_hq} & out & 4/0/0 & 4/0/0 & 2.7 & 2.5 \\
2 & \code{mpv\_gpunext} & out & 4/0/0 & 4/0/0 & 0.4 & 0.4 \\
2 & \code{mpv\_gpunext\_hq} & out & 4/0/0 & 4/0/0 & 10.4 & 10.3 \\
2 & \code{vkmark\_clear} & in & 9/0/0 & 9/0/0 & 0.0 & 0.0 \\
2 & \code{vkmark\_cube} & in & 4/0/0 & 4/0/0 & 2.0 & 1.9 \\
2 & \code{vkmark\_desktop} & in & 7/0/0 & 7/0/0 & 1.4 & 1.2 \\
2 & \code{vkmark\_effect2d} & in & 3/0/0 & 3/0/0 & 1.4 & 1.5 \\
2 & \code{vkmark\_shading} & in & 5/0/0 & 5/0/0 & 4.2 & 3.7 \\
2 & \code{vkmark\_texture} & in & 5/0/0 & 5/0/0 & 1.1 & 1.2 \\
2 & \code{vkmark\_vertex} & in & 5/0/0 & 5/0/0 & 1.1 & 1.2 \\
\bottomrule
\end{tabular}
\end{table*}


\bibliographystyle{ACM-Reference-Format}
\bibliography{refs}

\end{document}
